\documentclass[11pt]{article}

    \usepackage[breakable]{tcolorbox}
    \usepackage{parskip} % Stop auto-indenting (to mimic markdown behaviour)
    
    \usepackage{iftex}
    \ifPDFTeX
    	\usepackage[T1]{fontenc}
    	\usepackage{mathpazo}
    \else
    	\usepackage{fontspec}
    \fi

    % Basic figure setup, for now with no caption control since it's done
    % automatically by Pandoc (which extracts ![](path) syntax from Markdown).
    \usepackage{graphicx}
    % Maintain compatibility with old templates. Remove in nbconvert 6.0
    \let\Oldincludegraphics\includegraphics
    % Ensure that by default, figures have no caption (until we provide a
    % proper Figure object with a Caption API and a way to capture that
    % in the conversion process - todo).
    \usepackage{caption}
    \DeclareCaptionFormat{nocaption}{}
    \captionsetup{format=nocaption,aboveskip=0pt,belowskip=0pt}

    \usepackage[Export]{adjustbox} % Used to constrain images to a maximum size
    \adjustboxset{max size={0.9\linewidth}{0.9\paperheight}}
    \usepackage{float}
    \floatplacement{figure}{H} % forces figures to be placed at the correct location
    \usepackage{xcolor} % Allow colors to be defined
    \usepackage{enumerate} % Needed for markdown enumerations to work
    \usepackage{geometry} % Used to adjust the document margins
    \usepackage{amsmath} % Equations
    \usepackage{amssymb} % Equations
    \usepackage{textcomp} % defines textquotesingle
    % Hack from http://tex.stackexchange.com/a/47451/13684:
    \AtBeginDocument{%
        \def\PYZsq{\textquotesingle}% Upright quotes in Pygmentized code
    }
    \usepackage{upquote} % Upright quotes for verbatim code
    \usepackage{eurosym} % defines \euro
    \usepackage[mathletters]{ucs} % Extended unicode (utf-8) support
    \usepackage{fancyvrb} % verbatim replacement that allows latex
    \usepackage{grffile} % extends the file name processing of package graphics 
                         % to support a larger range
    \makeatletter % fix for grffile with XeLaTeX
    \def\Gread@@xetex#1{%
      \IfFileExists{"\Gin@base".bb}%
      {\Gread@eps{\Gin@base.bb}}%
      {\Gread@@xetex@aux#1}%
    }
    \makeatother

    % The hyperref package gives us a pdf with properly built
    % internal navigation ('pdf bookmarks' for the table of contents,
    % internal cross-reference links, web links for URLs, etc.)
    \usepackage{hyperref}
    % The default LaTeX title has an obnoxious amount of whitespace. By default,
    % titling removes some of it. It also provides customization options.
    \usepackage{titling}
    \usepackage{longtable} % longtable support required by pandoc >1.10
    \usepackage{booktabs}  % table support for pandoc > 1.12.2
    \usepackage[inline]{enumitem} % IRkernel/repr support (it uses the enumerate* environment)
    \usepackage[normalem]{ulem} % ulem is needed to support strikethroughs (\sout)
                                % normalem makes italics be italics, not underlines
    \usepackage{mathrsfs}
    

    
    % Colors for the hyperref package
    \definecolor{urlcolor}{rgb}{0,.145,.698}
    \definecolor{linkcolor}{rgb}{.71,0.21,0.01}
    \definecolor{citecolor}{rgb}{.12,.54,.11}

    % ANSI colors
    \definecolor{ansi-black}{HTML}{3E424D}
    \definecolor{ansi-black-intense}{HTML}{282C36}
    \definecolor{ansi-red}{HTML}{E75C58}
    \definecolor{ansi-red-intense}{HTML}{B22B31}
    \definecolor{ansi-green}{HTML}{00A250}
    \definecolor{ansi-green-intense}{HTML}{007427}
    \definecolor{ansi-yellow}{HTML}{DDB62B}
    \definecolor{ansi-yellow-intense}{HTML}{B27D12}
    \definecolor{ansi-blue}{HTML}{208FFB}
    \definecolor{ansi-blue-intense}{HTML}{0065CA}
    \definecolor{ansi-magenta}{HTML}{D160C4}
    \definecolor{ansi-magenta-intense}{HTML}{A03196}
    \definecolor{ansi-cyan}{HTML}{60C6C8}
    \definecolor{ansi-cyan-intense}{HTML}{258F8F}
    \definecolor{ansi-white}{HTML}{C5C1B4}
    \definecolor{ansi-white-intense}{HTML}{A1A6B2}
    \definecolor{ansi-default-inverse-fg}{HTML}{FFFFFF}
    \definecolor{ansi-default-inverse-bg}{HTML}{000000}

    % commands and environments needed by pandoc snippets
    % extracted from the output of `pandoc -s`
    \providecommand{\tightlist}{%
      \setlength{\itemsep}{0pt}\setlength{\parskip}{0pt}}
    \DefineVerbatimEnvironment{Highlighting}{Verbatim}{commandchars=\\\{\}}
    % Add ',fontsize=\small' for more characters per line
    \newenvironment{Shaded}{}{}
    \newcommand{\KeywordTok}[1]{\textcolor[rgb]{0.00,0.44,0.13}{\textbf{{#1}}}}
    \newcommand{\DataTypeTok}[1]{\textcolor[rgb]{0.56,0.13,0.00}{{#1}}}
    \newcommand{\DecValTok}[1]{\textcolor[rgb]{0.25,0.63,0.44}{{#1}}}
    \newcommand{\BaseNTok}[1]{\textcolor[rgb]{0.25,0.63,0.44}{{#1}}}
    \newcommand{\FloatTok}[1]{\textcolor[rgb]{0.25,0.63,0.44}{{#1}}}
    \newcommand{\CharTok}[1]{\textcolor[rgb]{0.25,0.44,0.63}{{#1}}}
    \newcommand{\StringTok}[1]{\textcolor[rgb]{0.25,0.44,0.63}{{#1}}}
    \newcommand{\CommentTok}[1]{\textcolor[rgb]{0.38,0.63,0.69}{\textit{{#1}}}}
    \newcommand{\OtherTok}[1]{\textcolor[rgb]{0.00,0.44,0.13}{{#1}}}
    \newcommand{\AlertTok}[1]{\textcolor[rgb]{1.00,0.00,0.00}{\textbf{{#1}}}}
    \newcommand{\FunctionTok}[1]{\textcolor[rgb]{0.02,0.16,0.49}{{#1}}}
    \newcommand{\RegionMarkerTok}[1]{{#1}}
    \newcommand{\ErrorTok}[1]{\textcolor[rgb]{1.00,0.00,0.00}{\textbf{{#1}}}}
    \newcommand{\NormalTok}[1]{{#1}}
    
    % Additional commands for more recent versions of Pandoc
    \newcommand{\ConstantTok}[1]{\textcolor[rgb]{0.53,0.00,0.00}{{#1}}}
    \newcommand{\SpecialCharTok}[1]{\textcolor[rgb]{0.25,0.44,0.63}{{#1}}}
    \newcommand{\VerbatimStringTok}[1]{\textcolor[rgb]{0.25,0.44,0.63}{{#1}}}
    \newcommand{\SpecialStringTok}[1]{\textcolor[rgb]{0.73,0.40,0.53}{{#1}}}
    \newcommand{\ImportTok}[1]{{#1}}
    \newcommand{\DocumentationTok}[1]{\textcolor[rgb]{0.73,0.13,0.13}{\textit{{#1}}}}
    \newcommand{\AnnotationTok}[1]{\textcolor[rgb]{0.38,0.63,0.69}{\textbf{\textit{{#1}}}}}
    \newcommand{\CommentVarTok}[1]{\textcolor[rgb]{0.38,0.63,0.69}{\textbf{\textit{{#1}}}}}
    \newcommand{\VariableTok}[1]{\textcolor[rgb]{0.10,0.09,0.49}{{#1}}}
    \newcommand{\ControlFlowTok}[1]{\textcolor[rgb]{0.00,0.44,0.13}{\textbf{{#1}}}}
    \newcommand{\OperatorTok}[1]{\textcolor[rgb]{0.40,0.40,0.40}{{#1}}}
    \newcommand{\BuiltInTok}[1]{{#1}}
    \newcommand{\ExtensionTok}[1]{{#1}}
    \newcommand{\PreprocessorTok}[1]{\textcolor[rgb]{0.74,0.48,0.00}{{#1}}}
    \newcommand{\AttributeTok}[1]{\textcolor[rgb]{0.49,0.56,0.16}{{#1}}}
    \newcommand{\InformationTok}[1]{\textcolor[rgb]{0.38,0.63,0.69}{\textbf{\textit{{#1}}}}}
    \newcommand{\WarningTok}[1]{\textcolor[rgb]{0.38,0.63,0.69}{\textbf{\textit{{#1}}}}}
    
    
    % Define a nice break command that doesn't care if a line doesn't already
    % exist.
    \def\br{\hspace*{\fill} \\* }
    % Math Jax compatibility definitions
    \def\gt{>}
    \def\lt{<}
    \let\Oldtex\TeX
    \let\Oldlatex\LaTeX
    \renewcommand{\TeX}{\textrm{\Oldtex}}
    \renewcommand{\LaTeX}{\textrm{\Oldlatex}}
    % Document parameters
    % Document title
    \title{Dinosaurus\_Island\_Character\_level\_language\_model}
    
    
    
    
    
% Pygments definitions
\makeatletter
\def\PY@reset{\let\PY@it=\relax \let\PY@bf=\relax%
    \let\PY@ul=\relax \let\PY@tc=\relax%
    \let\PY@bc=\relax \let\PY@ff=\relax}
\def\PY@tok#1{\csname PY@tok@#1\endcsname}
\def\PY@toks#1+{\ifx\relax#1\empty\else%
    \PY@tok{#1}\expandafter\PY@toks\fi}
\def\PY@do#1{\PY@bc{\PY@tc{\PY@ul{%
    \PY@it{\PY@bf{\PY@ff{#1}}}}}}}
\def\PY#1#2{\PY@reset\PY@toks#1+\relax+\PY@do{#2}}

\expandafter\def\csname PY@tok@w\endcsname{\def\PY@tc##1{\textcolor[rgb]{0.73,0.73,0.73}{##1}}}
\expandafter\def\csname PY@tok@c\endcsname{\let\PY@it=\textit\def\PY@tc##1{\textcolor[rgb]{0.25,0.50,0.50}{##1}}}
\expandafter\def\csname PY@tok@cp\endcsname{\def\PY@tc##1{\textcolor[rgb]{0.74,0.48,0.00}{##1}}}
\expandafter\def\csname PY@tok@k\endcsname{\let\PY@bf=\textbf\def\PY@tc##1{\textcolor[rgb]{0.00,0.50,0.00}{##1}}}
\expandafter\def\csname PY@tok@kp\endcsname{\def\PY@tc##1{\textcolor[rgb]{0.00,0.50,0.00}{##1}}}
\expandafter\def\csname PY@tok@kt\endcsname{\def\PY@tc##1{\textcolor[rgb]{0.69,0.00,0.25}{##1}}}
\expandafter\def\csname PY@tok@o\endcsname{\def\PY@tc##1{\textcolor[rgb]{0.40,0.40,0.40}{##1}}}
\expandafter\def\csname PY@tok@ow\endcsname{\let\PY@bf=\textbf\def\PY@tc##1{\textcolor[rgb]{0.67,0.13,1.00}{##1}}}
\expandafter\def\csname PY@tok@nb\endcsname{\def\PY@tc##1{\textcolor[rgb]{0.00,0.50,0.00}{##1}}}
\expandafter\def\csname PY@tok@nf\endcsname{\def\PY@tc##1{\textcolor[rgb]{0.00,0.00,1.00}{##1}}}
\expandafter\def\csname PY@tok@nc\endcsname{\let\PY@bf=\textbf\def\PY@tc##1{\textcolor[rgb]{0.00,0.00,1.00}{##1}}}
\expandafter\def\csname PY@tok@nn\endcsname{\let\PY@bf=\textbf\def\PY@tc##1{\textcolor[rgb]{0.00,0.00,1.00}{##1}}}
\expandafter\def\csname PY@tok@ne\endcsname{\let\PY@bf=\textbf\def\PY@tc##1{\textcolor[rgb]{0.82,0.25,0.23}{##1}}}
\expandafter\def\csname PY@tok@nv\endcsname{\def\PY@tc##1{\textcolor[rgb]{0.10,0.09,0.49}{##1}}}
\expandafter\def\csname PY@tok@no\endcsname{\def\PY@tc##1{\textcolor[rgb]{0.53,0.00,0.00}{##1}}}
\expandafter\def\csname PY@tok@nl\endcsname{\def\PY@tc##1{\textcolor[rgb]{0.63,0.63,0.00}{##1}}}
\expandafter\def\csname PY@tok@ni\endcsname{\let\PY@bf=\textbf\def\PY@tc##1{\textcolor[rgb]{0.60,0.60,0.60}{##1}}}
\expandafter\def\csname PY@tok@na\endcsname{\def\PY@tc##1{\textcolor[rgb]{0.49,0.56,0.16}{##1}}}
\expandafter\def\csname PY@tok@nt\endcsname{\let\PY@bf=\textbf\def\PY@tc##1{\textcolor[rgb]{0.00,0.50,0.00}{##1}}}
\expandafter\def\csname PY@tok@nd\endcsname{\def\PY@tc##1{\textcolor[rgb]{0.67,0.13,1.00}{##1}}}
\expandafter\def\csname PY@tok@s\endcsname{\def\PY@tc##1{\textcolor[rgb]{0.73,0.13,0.13}{##1}}}
\expandafter\def\csname PY@tok@sd\endcsname{\let\PY@it=\textit\def\PY@tc##1{\textcolor[rgb]{0.73,0.13,0.13}{##1}}}
\expandafter\def\csname PY@tok@si\endcsname{\let\PY@bf=\textbf\def\PY@tc##1{\textcolor[rgb]{0.73,0.40,0.53}{##1}}}
\expandafter\def\csname PY@tok@se\endcsname{\let\PY@bf=\textbf\def\PY@tc##1{\textcolor[rgb]{0.73,0.40,0.13}{##1}}}
\expandafter\def\csname PY@tok@sr\endcsname{\def\PY@tc##1{\textcolor[rgb]{0.73,0.40,0.53}{##1}}}
\expandafter\def\csname PY@tok@ss\endcsname{\def\PY@tc##1{\textcolor[rgb]{0.10,0.09,0.49}{##1}}}
\expandafter\def\csname PY@tok@sx\endcsname{\def\PY@tc##1{\textcolor[rgb]{0.00,0.50,0.00}{##1}}}
\expandafter\def\csname PY@tok@m\endcsname{\def\PY@tc##1{\textcolor[rgb]{0.40,0.40,0.40}{##1}}}
\expandafter\def\csname PY@tok@gh\endcsname{\let\PY@bf=\textbf\def\PY@tc##1{\textcolor[rgb]{0.00,0.00,0.50}{##1}}}
\expandafter\def\csname PY@tok@gu\endcsname{\let\PY@bf=\textbf\def\PY@tc##1{\textcolor[rgb]{0.50,0.00,0.50}{##1}}}
\expandafter\def\csname PY@tok@gd\endcsname{\def\PY@tc##1{\textcolor[rgb]{0.63,0.00,0.00}{##1}}}
\expandafter\def\csname PY@tok@gi\endcsname{\def\PY@tc##1{\textcolor[rgb]{0.00,0.63,0.00}{##1}}}
\expandafter\def\csname PY@tok@gr\endcsname{\def\PY@tc##1{\textcolor[rgb]{1.00,0.00,0.00}{##1}}}
\expandafter\def\csname PY@tok@ge\endcsname{\let\PY@it=\textit}
\expandafter\def\csname PY@tok@gs\endcsname{\let\PY@bf=\textbf}
\expandafter\def\csname PY@tok@gp\endcsname{\let\PY@bf=\textbf\def\PY@tc##1{\textcolor[rgb]{0.00,0.00,0.50}{##1}}}
\expandafter\def\csname PY@tok@go\endcsname{\def\PY@tc##1{\textcolor[rgb]{0.53,0.53,0.53}{##1}}}
\expandafter\def\csname PY@tok@gt\endcsname{\def\PY@tc##1{\textcolor[rgb]{0.00,0.27,0.87}{##1}}}
\expandafter\def\csname PY@tok@err\endcsname{\def\PY@bc##1{\setlength{\fboxsep}{0pt}\fcolorbox[rgb]{1.00,0.00,0.00}{1,1,1}{\strut ##1}}}
\expandafter\def\csname PY@tok@kc\endcsname{\let\PY@bf=\textbf\def\PY@tc##1{\textcolor[rgb]{0.00,0.50,0.00}{##1}}}
\expandafter\def\csname PY@tok@kd\endcsname{\let\PY@bf=\textbf\def\PY@tc##1{\textcolor[rgb]{0.00,0.50,0.00}{##1}}}
\expandafter\def\csname PY@tok@kn\endcsname{\let\PY@bf=\textbf\def\PY@tc##1{\textcolor[rgb]{0.00,0.50,0.00}{##1}}}
\expandafter\def\csname PY@tok@kr\endcsname{\let\PY@bf=\textbf\def\PY@tc##1{\textcolor[rgb]{0.00,0.50,0.00}{##1}}}
\expandafter\def\csname PY@tok@bp\endcsname{\def\PY@tc##1{\textcolor[rgb]{0.00,0.50,0.00}{##1}}}
\expandafter\def\csname PY@tok@fm\endcsname{\def\PY@tc##1{\textcolor[rgb]{0.00,0.00,1.00}{##1}}}
\expandafter\def\csname PY@tok@vc\endcsname{\def\PY@tc##1{\textcolor[rgb]{0.10,0.09,0.49}{##1}}}
\expandafter\def\csname PY@tok@vg\endcsname{\def\PY@tc##1{\textcolor[rgb]{0.10,0.09,0.49}{##1}}}
\expandafter\def\csname PY@tok@vi\endcsname{\def\PY@tc##1{\textcolor[rgb]{0.10,0.09,0.49}{##1}}}
\expandafter\def\csname PY@tok@vm\endcsname{\def\PY@tc##1{\textcolor[rgb]{0.10,0.09,0.49}{##1}}}
\expandafter\def\csname PY@tok@sa\endcsname{\def\PY@tc##1{\textcolor[rgb]{0.73,0.13,0.13}{##1}}}
\expandafter\def\csname PY@tok@sb\endcsname{\def\PY@tc##1{\textcolor[rgb]{0.73,0.13,0.13}{##1}}}
\expandafter\def\csname PY@tok@sc\endcsname{\def\PY@tc##1{\textcolor[rgb]{0.73,0.13,0.13}{##1}}}
\expandafter\def\csname PY@tok@dl\endcsname{\def\PY@tc##1{\textcolor[rgb]{0.73,0.13,0.13}{##1}}}
\expandafter\def\csname PY@tok@s2\endcsname{\def\PY@tc##1{\textcolor[rgb]{0.73,0.13,0.13}{##1}}}
\expandafter\def\csname PY@tok@sh\endcsname{\def\PY@tc##1{\textcolor[rgb]{0.73,0.13,0.13}{##1}}}
\expandafter\def\csname PY@tok@s1\endcsname{\def\PY@tc##1{\textcolor[rgb]{0.73,0.13,0.13}{##1}}}
\expandafter\def\csname PY@tok@mb\endcsname{\def\PY@tc##1{\textcolor[rgb]{0.40,0.40,0.40}{##1}}}
\expandafter\def\csname PY@tok@mf\endcsname{\def\PY@tc##1{\textcolor[rgb]{0.40,0.40,0.40}{##1}}}
\expandafter\def\csname PY@tok@mh\endcsname{\def\PY@tc##1{\textcolor[rgb]{0.40,0.40,0.40}{##1}}}
\expandafter\def\csname PY@tok@mi\endcsname{\def\PY@tc##1{\textcolor[rgb]{0.40,0.40,0.40}{##1}}}
\expandafter\def\csname PY@tok@il\endcsname{\def\PY@tc##1{\textcolor[rgb]{0.40,0.40,0.40}{##1}}}
\expandafter\def\csname PY@tok@mo\endcsname{\def\PY@tc##1{\textcolor[rgb]{0.40,0.40,0.40}{##1}}}
\expandafter\def\csname PY@tok@ch\endcsname{\let\PY@it=\textit\def\PY@tc##1{\textcolor[rgb]{0.25,0.50,0.50}{##1}}}
\expandafter\def\csname PY@tok@cm\endcsname{\let\PY@it=\textit\def\PY@tc##1{\textcolor[rgb]{0.25,0.50,0.50}{##1}}}
\expandafter\def\csname PY@tok@cpf\endcsname{\let\PY@it=\textit\def\PY@tc##1{\textcolor[rgb]{0.25,0.50,0.50}{##1}}}
\expandafter\def\csname PY@tok@c1\endcsname{\let\PY@it=\textit\def\PY@tc##1{\textcolor[rgb]{0.25,0.50,0.50}{##1}}}
\expandafter\def\csname PY@tok@cs\endcsname{\let\PY@it=\textit\def\PY@tc##1{\textcolor[rgb]{0.25,0.50,0.50}{##1}}}

\def\PYZbs{\char`\\}
\def\PYZus{\char`\_}
\def\PYZob{\char`\{}
\def\PYZcb{\char`\}}
\def\PYZca{\char`\^}
\def\PYZam{\char`\&}
\def\PYZlt{\char`\<}
\def\PYZgt{\char`\>}
\def\PYZsh{\char`\#}
\def\PYZpc{\char`\%}
\def\PYZdl{\char`\$}
\def\PYZhy{\char`\-}
\def\PYZsq{\char`\'}
\def\PYZdq{\char`\"}
\def\PYZti{\char`\~}
% for compatibility with earlier versions
\def\PYZat{@}
\def\PYZlb{[}
\def\PYZrb{]}
\makeatother


    % For linebreaks inside Verbatim environment from package fancyvrb. 
    \makeatletter
        \newbox\Wrappedcontinuationbox 
        \newbox\Wrappedvisiblespacebox 
        \newcommand*\Wrappedvisiblespace {\textcolor{red}{\textvisiblespace}} 
        \newcommand*\Wrappedcontinuationsymbol {\textcolor{red}{\llap{\tiny$\m@th\hookrightarrow$}}} 
        \newcommand*\Wrappedcontinuationindent {3ex } 
        \newcommand*\Wrappedafterbreak {\kern\Wrappedcontinuationindent\copy\Wrappedcontinuationbox} 
        % Take advantage of the already applied Pygments mark-up to insert 
        % potential linebreaks for TeX processing. 
        %        {, <, #, %, $, ' and ": go to next line. 
        %        _, }, ^, &, >, - and ~: stay at end of broken line. 
        % Use of \textquotesingle for straight quote. 
        \newcommand*\Wrappedbreaksatspecials {% 
            \def\PYGZus{\discretionary{\char`\_}{\Wrappedafterbreak}{\char`\_}}% 
            \def\PYGZob{\discretionary{}{\Wrappedafterbreak\char`\{}{\char`\{}}% 
            \def\PYGZcb{\discretionary{\char`\}}{\Wrappedafterbreak}{\char`\}}}% 
            \def\PYGZca{\discretionary{\char`\^}{\Wrappedafterbreak}{\char`\^}}% 
            \def\PYGZam{\discretionary{\char`\&}{\Wrappedafterbreak}{\char`\&}}% 
            \def\PYGZlt{\discretionary{}{\Wrappedafterbreak\char`\<}{\char`\<}}% 
            \def\PYGZgt{\discretionary{\char`\>}{\Wrappedafterbreak}{\char`\>}}% 
            \def\PYGZsh{\discretionary{}{\Wrappedafterbreak\char`\#}{\char`\#}}% 
            \def\PYGZpc{\discretionary{}{\Wrappedafterbreak\char`\%}{\char`\%}}% 
            \def\PYGZdl{\discretionary{}{\Wrappedafterbreak\char`\$}{\char`\$}}% 
            \def\PYGZhy{\discretionary{\char`\-}{\Wrappedafterbreak}{\char`\-}}% 
            \def\PYGZsq{\discretionary{}{\Wrappedafterbreak\textquotesingle}{\textquotesingle}}% 
            \def\PYGZdq{\discretionary{}{\Wrappedafterbreak\char`\"}{\char`\"}}% 
            \def\PYGZti{\discretionary{\char`\~}{\Wrappedafterbreak}{\char`\~}}% 
        } 
        % Some characters . , ; ? ! / are not pygmentized. 
        % This macro makes them "active" and they will insert potential linebreaks 
        \newcommand*\Wrappedbreaksatpunct {% 
            \lccode`\~`\.\lowercase{\def~}{\discretionary{\hbox{\char`\.}}{\Wrappedafterbreak}{\hbox{\char`\.}}}% 
            \lccode`\~`\,\lowercase{\def~}{\discretionary{\hbox{\char`\,}}{\Wrappedafterbreak}{\hbox{\char`\,}}}% 
            \lccode`\~`\;\lowercase{\def~}{\discretionary{\hbox{\char`\;}}{\Wrappedafterbreak}{\hbox{\char`\;}}}% 
            \lccode`\~`\:\lowercase{\def~}{\discretionary{\hbox{\char`\:}}{\Wrappedafterbreak}{\hbox{\char`\:}}}% 
            \lccode`\~`\?\lowercase{\def~}{\discretionary{\hbox{\char`\?}}{\Wrappedafterbreak}{\hbox{\char`\?}}}% 
            \lccode`\~`\!\lowercase{\def~}{\discretionary{\hbox{\char`\!}}{\Wrappedafterbreak}{\hbox{\char`\!}}}% 
            \lccode`\~`\/\lowercase{\def~}{\discretionary{\hbox{\char`\/}}{\Wrappedafterbreak}{\hbox{\char`\/}}}% 
            \catcode`\.\active
            \catcode`\,\active 
            \catcode`\;\active
            \catcode`\:\active
            \catcode`\?\active
            \catcode`\!\active
            \catcode`\/\active 
            \lccode`\~`\~ 	
        }
    \makeatother

    \let\OriginalVerbatim=\Verbatim
    \makeatletter
    \renewcommand{\Verbatim}[1][1]{%
        %\parskip\z@skip
        \sbox\Wrappedcontinuationbox {\Wrappedcontinuationsymbol}%
        \sbox\Wrappedvisiblespacebox {\FV@SetupFont\Wrappedvisiblespace}%
        \def\FancyVerbFormatLine ##1{\hsize\linewidth
            \vtop{\raggedright\hyphenpenalty\z@\exhyphenpenalty\z@
                \doublehyphendemerits\z@\finalhyphendemerits\z@
                \strut ##1\strut}%
        }%
        % If the linebreak is at a space, the latter will be displayed as visible
        % space at end of first line, and a continuation symbol starts next line.
        % Stretch/shrink are however usually zero for typewriter font.
        \def\FV@Space {%
            \nobreak\hskip\z@ plus\fontdimen3\font minus\fontdimen4\font
            \discretionary{\copy\Wrappedvisiblespacebox}{\Wrappedafterbreak}
            {\kern\fontdimen2\font}%
        }%
        
        % Allow breaks at special characters using \PYG... macros.
        \Wrappedbreaksatspecials
        % Breaks at punctuation characters . , ; ? ! and / need catcode=\active 	
        \OriginalVerbatim[#1,codes*=\Wrappedbreaksatpunct]%
    }
    \makeatother

    % Exact colors from NB
    \definecolor{incolor}{HTML}{303F9F}
    \definecolor{outcolor}{HTML}{D84315}
    \definecolor{cellborder}{HTML}{CFCFCF}
    \definecolor{cellbackground}{HTML}{F7F7F7}
    
    % prompt
    \makeatletter
    \newcommand{\boxspacing}{\kern\kvtcb@left@rule\kern\kvtcb@boxsep}
    \makeatother
    \newcommand{\prompt}[4]{
        \ttfamily\llap{{\color{#2}[#3]:\hspace{3pt}#4}}\vspace{-\baselineskip}
    }
    

    
    % Prevent overflowing lines due to hard-to-break entities
    \sloppy 
    % Setup hyperref package
    \hypersetup{
      breaklinks=true,  % so long urls are correctly broken across lines
      colorlinks=true,
      urlcolor=urlcolor,
      linkcolor=linkcolor,
      citecolor=citecolor,
      }
    % Slightly bigger margins than the latex defaults
    
    \geometry{verbose,tmargin=1in,bmargin=1in,lmargin=1in,rmargin=1in}
    
    

\begin{document}
    
    \maketitle
    
    

    
    \hypertarget{character-level-language-model---dinosaurus-island}{%
\section{Character level language model - Dinosaurus
Island}\label{character-level-language-model---dinosaurus-island}}

Welcome to Dinosaurus Island! 65 million years ago, dinosaurs existed,
and in this assignment, they have returned.

You are in charge of a special task: Leading biology researchers are
creating new breeds of dinosaurs and bringing them to life on earth, and
your job is to give names to these dinosaurs. If a dinosaur does not
like its name, it might go berserk, so choose wisely!

Luckily you're equipped with some deep learning now, and you will use it
to save the day! Your assistant has collected a list of all the dinosaur
names they could find, and compiled them into this
\href{dinos.txt}{dataset}. (Feel free to take a look by clicking the
previous link.) To create new dinosaur names, you will build a
character-level language model to generate new names. Your algorithm
will learn the different name patterns, and randomly generate new names.
Hopefully this algorithm will keep you and your team safe from the
dinosaurs' wrath!

By the time you complete this assignment, you'll be able to:

\begin{itemize}
\tightlist
\item
  Store text data for processing using an RNN
\item
  Build a character-level text generation model using an RNN
\item
  Sample novel sequences in an RNN
\item
  Explain the vanishing/exploding gradient problem in RNNs
\item
  Apply gradient clipping as a solution for exploding gradients
\end{itemize}

Begin by loading in some functions that are provided for you in
\texttt{rnn\_utils}. Specifically, you have access to functions such as
\texttt{rnn\_forward} and \texttt{rnn\_backward} which are equivalent to
those you've implemented in the previous assignment.

\hypertarget{important-note-on-submission-to-the-autograder}{%
\subsection{Important Note on Submission to the
AutoGrader}\label{important-note-on-submission-to-the-autograder}}

Before submitting your assignment to the AutoGrader, please make sure
you are not doing the following:

\begin{enumerate}
\def\labelenumi{\arabic{enumi}.}
\tightlist
\item
  You have not added any \emph{extra} \texttt{print} statement(s) in the
  assignment.
\item
  You have not added any \emph{extra} code cell(s) in the assignment.
\item
  You have not changed any of the function parameters.
\item
  You are not using any global variables inside your graded exercises.
  Unless specifically instructed to do so, please refrain from it and
  use the local variables instead.
\item
  You are not changing the assignment code where it is not required,
  like creating \emph{extra} variables.
\end{enumerate}

If you do any of the following, you will get something like,
\texttt{Grader\ Error:\ Grader\ feedback\ not\ found} (or similarly
unexpected) error upon submitting your assignment. Before asking for
help/debugging the errors in your assignment, check for these first. If
this is the case, and you don't remember the changes you have made, you
can get a fresh copy of the assignment by following these
\href{https://www.coursera.org/learn/nlp-sequence-models/supplement/qHIve/h-ow-to-refresh-your-workspace}{instructions}.

    \hypertarget{table-of-contents}{%
\subsection{Table of Contents}\label{table-of-contents}}

\begin{itemize}
\tightlist
\item
  Section \ref{0}
\item
  Section \ref{1}

  \begin{itemize}
  \tightlist
  \item
    Section \ref{1-1}
  \item
    Section \ref{1-2}
  \end{itemize}
\item
  Section \ref{2}

  \begin{itemize}
  \tightlist
  \item
    Section \ref{2-1}

    \begin{itemize}
    \tightlist
    \item
      Section \ref{ex-1}
    \end{itemize}
  \item
    Section \ref{2-2}

    \begin{itemize}
    \tightlist
    \item
      Section \ref{ex-2}
    \end{itemize}
  \end{itemize}
\item
  Section \ref{3}

  \begin{itemize}
  \tightlist
  \item
    Section \ref{3-1}

    \begin{itemize}
    \tightlist
    \item
      Section \ref{ex-3}
    \end{itemize}
  \item
    Section \ref{3-2}

    \begin{itemize}
    \tightlist
    \item
      Section \ref{ex-4}
    \end{itemize}
  \end{itemize}
\item
  Section \ref{4}
\item
  Section \ref{5}
\end{itemize}

    \#\# Packages

    \begin{tcolorbox}[breakable, size=fbox, boxrule=1pt, pad at break*=1mm,colback=cellbackground, colframe=cellborder]
\prompt{In}{incolor}{1}{\boxspacing}
\begin{Verbatim}[commandchars=\\\{\}]
\PY{c+c1}{\PYZsh{}\PYZsh{}\PYZsh{} v1.1}
\end{Verbatim}
\end{tcolorbox}

    \begin{tcolorbox}[breakable, size=fbox, boxrule=1pt, pad at break*=1mm,colback=cellbackground, colframe=cellborder]
\prompt{In}{incolor}{2}{\boxspacing}
\begin{Verbatim}[commandchars=\\\{\}]
\PY{k+kn}{import} \PY{n+nn}{numpy} \PY{k}{as} \PY{n+nn}{np}
\PY{k+kn}{from} \PY{n+nn}{utils} \PY{k+kn}{import} \PY{o}{*}
\PY{k+kn}{import} \PY{n+nn}{random}
\PY{k+kn}{import} \PY{n+nn}{pprint}
\PY{k+kn}{import} \PY{n+nn}{copy}
\end{Verbatim}
\end{tcolorbox}

    \#\# 1 - Problem Statement

\#\#\# 1.1 - Dataset and Preprocessing

Run the following cell to read the dataset of dinosaur names, create a
list of unique characters (such as a-z), and compute the dataset and
vocabulary size.

    \begin{tcolorbox}[breakable, size=fbox, boxrule=1pt, pad at break*=1mm,colback=cellbackground, colframe=cellborder]
\prompt{In}{incolor}{3}{\boxspacing}
\begin{Verbatim}[commandchars=\\\{\}]
\PY{n}{data} \PY{o}{=} \PY{n+nb}{open}\PY{p}{(}\PY{l+s+s1}{\PYZsq{}}\PY{l+s+s1}{dinos.txt}\PY{l+s+s1}{\PYZsq{}}\PY{p}{,} \PY{l+s+s1}{\PYZsq{}}\PY{l+s+s1}{r}\PY{l+s+s1}{\PYZsq{}}\PY{p}{)}\PY{o}{.}\PY{n}{read}\PY{p}{(}\PY{p}{)}
\PY{n}{data}\PY{o}{=} \PY{n}{data}\PY{o}{.}\PY{n}{lower}\PY{p}{(}\PY{p}{)}
\PY{n}{chars} \PY{o}{=} \PY{n+nb}{list}\PY{p}{(}\PY{n+nb}{set}\PY{p}{(}\PY{n}{data}\PY{p}{)}\PY{p}{)}
\PY{n}{data\PYZus{}size}\PY{p}{,} \PY{n}{vocab\PYZus{}size} \PY{o}{=} \PY{n+nb}{len}\PY{p}{(}\PY{n}{data}\PY{p}{)}\PY{p}{,} \PY{n+nb}{len}\PY{p}{(}\PY{n}{chars}\PY{p}{)}
\PY{n+nb}{print}\PY{p}{(}\PY{l+s+s1}{\PYZsq{}}\PY{l+s+s1}{There are }\PY{l+s+si}{\PYZpc{}d}\PY{l+s+s1}{ total characters and }\PY{l+s+si}{\PYZpc{}d}\PY{l+s+s1}{ unique characters in your data.}\PY{l+s+s1}{\PYZsq{}} \PY{o}{\PYZpc{}} \PY{p}{(}\PY{n}{data\PYZus{}size}\PY{p}{,} \PY{n}{vocab\PYZus{}size}\PY{p}{)}\PY{p}{)}
\end{Verbatim}
\end{tcolorbox}

    \begin{Verbatim}[commandchars=\\\{\}]
There are 19909 total characters and 27 unique characters in your data.
    \end{Verbatim}

    \begin{itemize}
\tightlist
\item
  The characters are a-z (26 characters) plus the ``\n'' (or newline
  character).
\item
  In this assignment, the newline character ``\n'' plays a role similar
  to the \texttt{\textless{}EOS\textgreater{}} (or ``End of sentence'')
  token discussed in lecture.

  \begin{itemize}
  \tightlist
  \item
    Here, ``\n'' indicates the end of the dinosaur name rather than the
    end of a sentence.
  \end{itemize}
\item
  \texttt{char\_to\_ix}: In the cell below, you'll create a Python
  dictionary (i.e., a hash table) to map each character to an index from
  0-26.
\item
  \texttt{ix\_to\_char}: Then, you'll create a second Python dictionary
  that maps each index back to the corresponding character.

  \begin{itemize}
  \tightlist
  \item
    This will help you figure out which index corresponds to which
    character in the probability distribution output of the softmax
    layer.
  \end{itemize}
\end{itemize}

    \begin{tcolorbox}[breakable, size=fbox, boxrule=1pt, pad at break*=1mm,colback=cellbackground, colframe=cellborder]
\prompt{In}{incolor}{4}{\boxspacing}
\begin{Verbatim}[commandchars=\\\{\}]
\PY{n}{chars} \PY{o}{=} \PY{n+nb}{sorted}\PY{p}{(}\PY{n}{chars}\PY{p}{)}
\PY{n+nb}{print}\PY{p}{(}\PY{n}{chars}\PY{p}{)}
\end{Verbatim}
\end{tcolorbox}

    \begin{Verbatim}[commandchars=\\\{\}]
['\textbackslash{}n', 'a', 'b', 'c', 'd', 'e', 'f', 'g', 'h', 'i', 'j', 'k', 'l', 'm', 'n',
'o', 'p', 'q', 'r', 's', 't', 'u', 'v', 'w', 'x', 'y', 'z']
    \end{Verbatim}

    \begin{tcolorbox}[breakable, size=fbox, boxrule=1pt, pad at break*=1mm,colback=cellbackground, colframe=cellborder]
\prompt{In}{incolor}{5}{\boxspacing}
\begin{Verbatim}[commandchars=\\\{\}]
\PY{n}{char\PYZus{}to\PYZus{}ix} \PY{o}{=} \PY{p}{\PYZob{}} \PY{n}{ch}\PY{p}{:}\PY{n}{i} \PY{k}{for} \PY{n}{i}\PY{p}{,}\PY{n}{ch} \PY{o+ow}{in} \PY{n+nb}{enumerate}\PY{p}{(}\PY{n}{chars}\PY{p}{)} \PY{p}{\PYZcb{}}
\PY{n}{ix\PYZus{}to\PYZus{}char} \PY{o}{=} \PY{p}{\PYZob{}} \PY{n}{i}\PY{p}{:}\PY{n}{ch} \PY{k}{for} \PY{n}{i}\PY{p}{,}\PY{n}{ch} \PY{o+ow}{in} \PY{n+nb}{enumerate}\PY{p}{(}\PY{n}{chars}\PY{p}{)} \PY{p}{\PYZcb{}}
\PY{n}{pp} \PY{o}{=} \PY{n}{pprint}\PY{o}{.}\PY{n}{PrettyPrinter}\PY{p}{(}\PY{n}{indent}\PY{o}{=}\PY{l+m+mi}{4}\PY{p}{)}
\PY{n}{pp}\PY{o}{.}\PY{n}{pprint}\PY{p}{(}\PY{n}{ix\PYZus{}to\PYZus{}char}\PY{p}{)}
\end{Verbatim}
\end{tcolorbox}

    \begin{Verbatim}[commandchars=\\\{\}]
\{   0: '\textbackslash{}n',
    1: 'a',
    2: 'b',
    3: 'c',
    4: 'd',
    5: 'e',
    6: 'f',
    7: 'g',
    8: 'h',
    9: 'i',
    10: 'j',
    11: 'k',
    12: 'l',
    13: 'm',
    14: 'n',
    15: 'o',
    16: 'p',
    17: 'q',
    18: 'r',
    19: 's',
    20: 't',
    21: 'u',
    22: 'v',
    23: 'w',
    24: 'x',
    25: 'y',
    26: 'z'\}
    \end{Verbatim}

    \#\#\# 1.2 - Overview of the Model

Your model will have the following structure:

\begin{itemize}
\tightlist
\item
  Initialize parameters
\item
  Run the optimization loop

  \begin{itemize}
  \tightlist
  \item
    Forward propagation to compute the loss function
  \item
    Backward propagation to compute the gradients with respect to the
    loss function
  \item
    Clip the gradients to avoid exploding gradients
  \item
    Using the gradients, update your parameters with the gradient
    descent update rule.
  \end{itemize}
\item
  Return the learned parameters
\end{itemize}

Figure 1: Recurrent Neural Network, similar to what you built in the
previous notebook ``Building a Recurrent Neural Network - Step by
Step.''

\begin{itemize}
\tightlist
\item
  At each time-step, the RNN tries to predict what the next character
  is, given the previous characters.
\item
  \(\mathbf{X} = (x^{\langle 1 \rangle}, x^{\langle 2 \rangle}, ..., x^{\langle T_x \rangle})\)
  is a list of characters from the training set.
\item
  \(\mathbf{Y} = (y^{\langle 1 \rangle}, y^{\langle 2 \rangle}, ..., y^{\langle T_x \rangle})\)
  is the same list of characters but shifted one character forward.
\item
  At every time-step \(t\),
  \(y^{\langle t \rangle} = x^{\langle t+1 \rangle}\). The prediction at
  time \(t\) is the same as the input at time \(t + 1\).
\end{itemize}

    \#\# 2 - Building Blocks of the Model

In this part, you will build two important blocks of the overall model:

\begin{enumerate}
\def\labelenumi{\arabic{enumi}.}
\tightlist
\item
  Gradient clipping: to avoid exploding gradients
\item
  Sampling: a technique used to generate characters
\end{enumerate}

You will then apply these two functions to build the model.

    \#\#\# 2.1 - Clipping the Gradients in the Optimization Loop

In this section you will implement the \texttt{clip} function that you
will call inside of your optimization loop.

\hypertarget{exploding-gradients}{%
\paragraph{Exploding gradients}\label{exploding-gradients}}

\begin{itemize}
\tightlist
\item
  When gradients are very large, they're called ``exploding
  gradients.''\\
\item
  Exploding gradients make the training process more difficult, because
  the updates may be so large that they ``overshoot'' the optimal values
  during back propagation.
\end{itemize}

Recall that your overall loop structure usually consists of: * forward
pass, * cost computation, * backward pass, * parameter update.

Before updating the parameters, you will perform gradient clipping to
make sure that your gradients are not ``exploding.''

\hypertarget{gradient-clipping}{%
\paragraph{Gradient clipping}\label{gradient-clipping}}

In the exercise below, you will implement a function \texttt{clip} that
takes in a dictionary of gradients and returns a clipped version of
gradients, if needed.

\begin{itemize}
\tightlist
\item
  There are different ways to clip gradients.
\item
  You will use a simple element-wise clipping procedure, in which every
  element of the gradient vector is clipped to fall between some range
  {[}-N, N{]}.
\item
  For example, if the N=10

  \begin{itemize}
  \tightlist
  \item
    The range is {[}-10, 10{]}
  \item
    If any component of the gradient vector is greater than 10, it is
    set to 10.
  \item
    If any component of the gradient vector is less than -10, it is set
    to -10.
  \item
    If any components are between -10 and 10, they keep their original
    values.
  \end{itemize}
\end{itemize}

Figure 2: Visualization of gradient descent with and without gradient
clipping, in a case where the network is running into ``exploding
gradient'' problems.

\#\#\# Exercise 1 - clip

Return the clipped gradients of your dictionary \texttt{gradients}.

\begin{itemize}
\tightlist
\item
  Your function takes in a maximum threshold and returns the clipped
  versions of the gradients.
\item
  You can check out
  \href{https://docs.scipy.org/doc/numpy-1.13.0/reference/generated/numpy.clip.html}{numpy.clip}
  for more info.

  \begin{itemize}
  \tightlist
  \item
    You will need to use the argument ``\texttt{out\ =\ ...}''.
  \item
    Using the ``\texttt{out}'' parameter allows you to update a variable
    ``in-place''.
  \item
    If you don't use ``\texttt{out}'' argument, the clipped variable is
    stored in the variable ``gradient'' but does not update the gradient
    variables \texttt{dWax}, \texttt{dWaa}, \texttt{dWya}, \texttt{db},
    \texttt{dby}.
  \end{itemize}
\end{itemize}

    \begin{tcolorbox}[breakable, size=fbox, boxrule=1pt, pad at break*=1mm,colback=cellbackground, colframe=cellborder]
\prompt{In}{incolor}{34}{\boxspacing}
\begin{Verbatim}[commandchars=\\\{\}]
\PY{c+c1}{\PYZsh{} UNQ\PYZus{}C1 (UNIQUE CELL IDENTIFIER, DO NOT EDIT)}
\PY{c+c1}{\PYZsh{}\PYZsh{}\PYZsh{} GRADED FUNCTION: clip}

\PY{k}{def} \PY{n+nf}{clip}\PY{p}{(}\PY{n}{gradients}\PY{p}{,} \PY{n}{maxValue}\PY{p}{)}\PY{p}{:}
    \PY{l+s+sd}{\PYZsq{}\PYZsq{}\PYZsq{}}
\PY{l+s+sd}{    Clips the gradients\PYZsq{} values between minimum and maximum.}
\PY{l+s+sd}{    }
\PY{l+s+sd}{    Arguments:}
\PY{l+s+sd}{    gradients \PYZhy{}\PYZhy{} a dictionary containing the gradients \PYZdq{}dWaa\PYZdq{}, \PYZdq{}dWax\PYZdq{}, \PYZdq{}dWya\PYZdq{}, \PYZdq{}db\PYZdq{}, \PYZdq{}dby\PYZdq{}}
\PY{l+s+sd}{    maxValue \PYZhy{}\PYZhy{} everything above this number is set to this number, and everything less than \PYZhy{}maxValue is set to \PYZhy{}maxValue}
\PY{l+s+sd}{    }
\PY{l+s+sd}{    Returns: }
\PY{l+s+sd}{    gradients \PYZhy{}\PYZhy{} a dictionary with the clipped gradients.}
\PY{l+s+sd}{    \PYZsq{}\PYZsq{}\PYZsq{}}
    \PY{n}{gradients} \PY{o}{=} \PY{n}{copy}\PY{o}{.}\PY{n}{deepcopy}\PY{p}{(}\PY{n}{gradients}\PY{p}{)}
    
    \PY{n}{dWaa}\PY{p}{,} \PY{n}{dWax}\PY{p}{,} \PY{n}{dWya}\PY{p}{,} \PY{n}{db}\PY{p}{,} \PY{n}{dby} \PY{o}{=} \PY{n}{gradients}\PY{p}{[}\PY{l+s+s1}{\PYZsq{}}\PY{l+s+s1}{dWaa}\PY{l+s+s1}{\PYZsq{}}\PY{p}{]}\PY{p}{,} \PY{n}{gradients}\PY{p}{[}\PY{l+s+s1}{\PYZsq{}}\PY{l+s+s1}{dWax}\PY{l+s+s1}{\PYZsq{}}\PY{p}{]}\PY{p}{,} \PY{n}{gradients}\PY{p}{[}\PY{l+s+s1}{\PYZsq{}}\PY{l+s+s1}{dWya}\PY{l+s+s1}{\PYZsq{}}\PY{p}{]}\PY{p}{,} \PY{n}{gradients}\PY{p}{[}\PY{l+s+s1}{\PYZsq{}}\PY{l+s+s1}{db}\PY{l+s+s1}{\PYZsq{}}\PY{p}{]}\PY{p}{,} \PY{n}{gradients}\PY{p}{[}\PY{l+s+s1}{\PYZsq{}}\PY{l+s+s1}{dby}\PY{l+s+s1}{\PYZsq{}}\PY{p}{]}
   
    \PY{c+c1}{\PYZsh{}\PYZsh{}\PYZsh{} START CODE HERE \PYZsh{}\PYZsh{}\PYZsh{}}
    \PY{c+c1}{\PYZsh{} Clip to mitigate exploding gradients, loop over [dWax, dWaa, dWya, db, dby]. (≈2 lines)}
    \PY{k}{for} \PY{n}{gradient} \PY{o+ow}{in} \PY{n}{gradients}\PY{p}{:}
        \PY{n}{np}\PY{o}{.}\PY{n}{clip}\PY{p}{(}\PY{n}{gradients}\PY{p}{[}\PY{n}{gradient}\PY{p}{]}\PY{p}{,} \PY{o}{\PYZhy{}}\PY{n}{maxValue}\PY{p}{,} \PY{n}{maxValue}\PY{p}{,} \PY{n}{out} \PY{o}{=} \PY{n}{gradients}\PY{p}{[}\PY{n}{gradient}\PY{p}{]}\PY{p}{)}
    \PY{c+c1}{\PYZsh{}\PYZsh{}\PYZsh{} END CODE HERE \PYZsh{}\PYZsh{}\PYZsh{}}
    
    \PY{n}{gradients} \PY{o}{=} \PY{p}{\PYZob{}}\PY{l+s+s2}{\PYZdq{}}\PY{l+s+s2}{dWaa}\PY{l+s+s2}{\PYZdq{}}\PY{p}{:} \PY{n}{dWaa}\PY{p}{,} \PY{l+s+s2}{\PYZdq{}}\PY{l+s+s2}{dWax}\PY{l+s+s2}{\PYZdq{}}\PY{p}{:} \PY{n}{dWax}\PY{p}{,} \PY{l+s+s2}{\PYZdq{}}\PY{l+s+s2}{dWya}\PY{l+s+s2}{\PYZdq{}}\PY{p}{:} \PY{n}{dWya}\PY{p}{,} \PY{l+s+s2}{\PYZdq{}}\PY{l+s+s2}{db}\PY{l+s+s2}{\PYZdq{}}\PY{p}{:} \PY{n}{db}\PY{p}{,} \PY{l+s+s2}{\PYZdq{}}\PY{l+s+s2}{dby}\PY{l+s+s2}{\PYZdq{}}\PY{p}{:} \PY{n}{dby}\PY{p}{\PYZcb{}}
    
    \PY{k}{return} \PY{n}{gradients}
\end{Verbatim}
\end{tcolorbox}

    \begin{tcolorbox}[breakable, size=fbox, boxrule=1pt, pad at break*=1mm,colback=cellbackground, colframe=cellborder]
\prompt{In}{incolor}{35}{\boxspacing}
\begin{Verbatim}[commandchars=\\\{\}]
\PY{c+c1}{\PYZsh{} Test with a max value of 10}
\PY{k}{def} \PY{n+nf}{clip\PYZus{}test}\PY{p}{(}\PY{n}{target}\PY{p}{,} \PY{n}{mValue}\PY{p}{)}\PY{p}{:}
    \PY{n+nb}{print}\PY{p}{(}\PY{l+s+sa}{f}\PY{l+s+s2}{\PYZdq{}}\PY{l+s+se}{\PYZbs{}n}\PY{l+s+s2}{Gradients for mValue=}\PY{l+s+si}{\PYZob{}}\PY{n}{mValue}\PY{l+s+si}{\PYZcb{}}\PY{l+s+s2}{\PYZdq{}}\PY{p}{)}
    \PY{n}{np}\PY{o}{.}\PY{n}{random}\PY{o}{.}\PY{n}{seed}\PY{p}{(}\PY{l+m+mi}{3}\PY{p}{)}
    \PY{n}{dWax} \PY{o}{=} \PY{n}{np}\PY{o}{.}\PY{n}{random}\PY{o}{.}\PY{n}{randn}\PY{p}{(}\PY{l+m+mi}{5}\PY{p}{,} \PY{l+m+mi}{3}\PY{p}{)} \PY{o}{*} \PY{l+m+mi}{10}
    \PY{n}{dWaa} \PY{o}{=} \PY{n}{np}\PY{o}{.}\PY{n}{random}\PY{o}{.}\PY{n}{randn}\PY{p}{(}\PY{l+m+mi}{5}\PY{p}{,} \PY{l+m+mi}{5}\PY{p}{)} \PY{o}{*} \PY{l+m+mi}{10}
    \PY{n}{dWya} \PY{o}{=} \PY{n}{np}\PY{o}{.}\PY{n}{random}\PY{o}{.}\PY{n}{randn}\PY{p}{(}\PY{l+m+mi}{2}\PY{p}{,} \PY{l+m+mi}{5}\PY{p}{)} \PY{o}{*} \PY{l+m+mi}{10}
    \PY{n}{db} \PY{o}{=} \PY{n}{np}\PY{o}{.}\PY{n}{random}\PY{o}{.}\PY{n}{randn}\PY{p}{(}\PY{l+m+mi}{5}\PY{p}{,} \PY{l+m+mi}{1}\PY{p}{)} \PY{o}{*} \PY{l+m+mi}{10}
    \PY{n}{dby} \PY{o}{=} \PY{n}{np}\PY{o}{.}\PY{n}{random}\PY{o}{.}\PY{n}{randn}\PY{p}{(}\PY{l+m+mi}{2}\PY{p}{,} \PY{l+m+mi}{1}\PY{p}{)} \PY{o}{*} \PY{l+m+mi}{10}
    \PY{n}{gradients} \PY{o}{=} \PY{p}{\PYZob{}}\PY{l+s+s2}{\PYZdq{}}\PY{l+s+s2}{dWax}\PY{l+s+s2}{\PYZdq{}}\PY{p}{:} \PY{n}{dWax}\PY{p}{,} \PY{l+s+s2}{\PYZdq{}}\PY{l+s+s2}{dWaa}\PY{l+s+s2}{\PYZdq{}}\PY{p}{:} \PY{n}{dWaa}\PY{p}{,} \PY{l+s+s2}{\PYZdq{}}\PY{l+s+s2}{dWya}\PY{l+s+s2}{\PYZdq{}}\PY{p}{:} \PY{n}{dWya}\PY{p}{,} \PY{l+s+s2}{\PYZdq{}}\PY{l+s+s2}{db}\PY{l+s+s2}{\PYZdq{}}\PY{p}{:} \PY{n}{db}\PY{p}{,} \PY{l+s+s2}{\PYZdq{}}\PY{l+s+s2}{dby}\PY{l+s+s2}{\PYZdq{}}\PY{p}{:} \PY{n}{dby}\PY{p}{\PYZcb{}}

    \PY{n}{gradients2} \PY{o}{=} \PY{n}{target}\PY{p}{(}\PY{n}{gradients}\PY{p}{,} \PY{n}{mValue}\PY{p}{)}
    \PY{n+nb}{print}\PY{p}{(}\PY{l+s+s2}{\PYZdq{}}\PY{l+s+s2}{gradients[}\PY{l+s+se}{\PYZbs{}\PYZdq{}}\PY{l+s+s2}{dWaa}\PY{l+s+se}{\PYZbs{}\PYZdq{}}\PY{l+s+s2}{][1][2] =}\PY{l+s+s2}{\PYZdq{}}\PY{p}{,} \PY{n}{gradients2}\PY{p}{[}\PY{l+s+s2}{\PYZdq{}}\PY{l+s+s2}{dWaa}\PY{l+s+s2}{\PYZdq{}}\PY{p}{]}\PY{p}{[}\PY{l+m+mi}{1}\PY{p}{]}\PY{p}{[}\PY{l+m+mi}{2}\PY{p}{]}\PY{p}{)}
    \PY{n+nb}{print}\PY{p}{(}\PY{l+s+s2}{\PYZdq{}}\PY{l+s+s2}{gradients[}\PY{l+s+se}{\PYZbs{}\PYZdq{}}\PY{l+s+s2}{dWax}\PY{l+s+se}{\PYZbs{}\PYZdq{}}\PY{l+s+s2}{][3][1] =}\PY{l+s+s2}{\PYZdq{}}\PY{p}{,} \PY{n}{gradients2}\PY{p}{[}\PY{l+s+s2}{\PYZdq{}}\PY{l+s+s2}{dWax}\PY{l+s+s2}{\PYZdq{}}\PY{p}{]}\PY{p}{[}\PY{l+m+mi}{3}\PY{p}{]}\PY{p}{[}\PY{l+m+mi}{1}\PY{p}{]}\PY{p}{)}
    \PY{n+nb}{print}\PY{p}{(}\PY{l+s+s2}{\PYZdq{}}\PY{l+s+s2}{gradients[}\PY{l+s+se}{\PYZbs{}\PYZdq{}}\PY{l+s+s2}{dWya}\PY{l+s+se}{\PYZbs{}\PYZdq{}}\PY{l+s+s2}{][1][2] =}\PY{l+s+s2}{\PYZdq{}}\PY{p}{,} \PY{n}{gradients2}\PY{p}{[}\PY{l+s+s2}{\PYZdq{}}\PY{l+s+s2}{dWya}\PY{l+s+s2}{\PYZdq{}}\PY{p}{]}\PY{p}{[}\PY{l+m+mi}{1}\PY{p}{]}\PY{p}{[}\PY{l+m+mi}{2}\PY{p}{]}\PY{p}{)}
    \PY{n+nb}{print}\PY{p}{(}\PY{l+s+s2}{\PYZdq{}}\PY{l+s+s2}{gradients[}\PY{l+s+se}{\PYZbs{}\PYZdq{}}\PY{l+s+s2}{db}\PY{l+s+se}{\PYZbs{}\PYZdq{}}\PY{l+s+s2}{][4] =}\PY{l+s+s2}{\PYZdq{}}\PY{p}{,} \PY{n}{gradients2}\PY{p}{[}\PY{l+s+s2}{\PYZdq{}}\PY{l+s+s2}{db}\PY{l+s+s2}{\PYZdq{}}\PY{p}{]}\PY{p}{[}\PY{l+m+mi}{4}\PY{p}{]}\PY{p}{)}
    \PY{n+nb}{print}\PY{p}{(}\PY{l+s+s2}{\PYZdq{}}\PY{l+s+s2}{gradients[}\PY{l+s+se}{\PYZbs{}\PYZdq{}}\PY{l+s+s2}{dby}\PY{l+s+se}{\PYZbs{}\PYZdq{}}\PY{l+s+s2}{][1] =}\PY{l+s+s2}{\PYZdq{}}\PY{p}{,} \PY{n}{gradients2}\PY{p}{[}\PY{l+s+s2}{\PYZdq{}}\PY{l+s+s2}{dby}\PY{l+s+s2}{\PYZdq{}}\PY{p}{]}\PY{p}{[}\PY{l+m+mi}{1}\PY{p}{]}\PY{p}{)}
    
    \PY{k}{for} \PY{n}{grad} \PY{o+ow}{in} \PY{n}{gradients2}\PY{o}{.}\PY{n}{keys}\PY{p}{(}\PY{p}{)}\PY{p}{:}
        \PY{n}{valuei} \PY{o}{=} \PY{n}{gradients}\PY{p}{[}\PY{n}{grad}\PY{p}{]}
        \PY{n}{valuef} \PY{o}{=} \PY{n}{gradients2}\PY{p}{[}\PY{n}{grad}\PY{p}{]}
        \PY{n}{mink} \PY{o}{=} \PY{n}{np}\PY{o}{.}\PY{n}{min}\PY{p}{(}\PY{n}{valuef}\PY{p}{)}
        \PY{n}{maxk} \PY{o}{=} \PY{n}{np}\PY{o}{.}\PY{n}{max}\PY{p}{(}\PY{n}{valuef}\PY{p}{)}
        \PY{k}{assert} \PY{n}{mink} \PY{o}{\PYZgt{}}\PY{o}{=} \PY{o}{\PYZhy{}}\PY{n+nb}{abs}\PY{p}{(}\PY{n}{mValue}\PY{p}{)}\PY{p}{,} \PY{l+s+sa}{f}\PY{l+s+s2}{\PYZdq{}}\PY{l+s+s2}{Problem with }\PY{l+s+si}{\PYZob{}}\PY{n}{grad}\PY{l+s+si}{\PYZcb{}}\PY{l+s+s2}{. Set a\PYZus{}min to \PYZhy{}mValue in the np.clip call}\PY{l+s+s2}{\PYZdq{}}
        \PY{k}{assert} \PY{n}{maxk} \PY{o}{\PYZlt{}}\PY{o}{=} \PY{n+nb}{abs}\PY{p}{(}\PY{n}{mValue}\PY{p}{)}\PY{p}{,} \PY{l+s+sa}{f}\PY{l+s+s2}{\PYZdq{}}\PY{l+s+s2}{Problem with }\PY{l+s+si}{\PYZob{}}\PY{n}{grad}\PY{l+s+si}{\PYZcb{}}\PY{l+s+s2}{.Set a\PYZus{}max to mValue in the np.clip call}\PY{l+s+s2}{\PYZdq{}}
        \PY{n}{index\PYZus{}not\PYZus{}clipped} \PY{o}{=} \PY{n}{np}\PY{o}{.}\PY{n}{logical\PYZus{}and}\PY{p}{(}\PY{n}{valuei} \PY{o}{\PYZlt{}}\PY{o}{=} \PY{n}{mValue}\PY{p}{,} \PY{n}{valuei} \PY{o}{\PYZgt{}}\PY{o}{=} \PY{o}{\PYZhy{}}\PY{n}{mValue}\PY{p}{)}
        \PY{k}{assert} \PY{n}{np}\PY{o}{.}\PY{n}{all}\PY{p}{(}\PY{n}{valuei}\PY{p}{[}\PY{n}{index\PYZus{}not\PYZus{}clipped}\PY{p}{]} \PY{o}{==} \PY{n}{valuef}\PY{p}{[}\PY{n}{index\PYZus{}not\PYZus{}clipped}\PY{p}{]}\PY{p}{)}\PY{p}{,} \PY{l+s+sa}{f}\PY{l+s+s2}{\PYZdq{}}\PY{l+s+s2}{ Problem with }\PY{l+s+si}{\PYZob{}}\PY{n}{grad}\PY{l+s+si}{\PYZcb{}}\PY{l+s+s2}{. Some values that should not have changed, changed during the clipping process.}\PY{l+s+s2}{\PYZdq{}}
    
    \PY{n+nb}{print}\PY{p}{(}\PY{l+s+s2}{\PYZdq{}}\PY{l+s+se}{\PYZbs{}033}\PY{l+s+s2}{[92mAll tests passed!}\PY{l+s+se}{\PYZbs{}x1b}\PY{l+s+s2}{[0m}\PY{l+s+s2}{\PYZdq{}}\PY{p}{)}
    
\PY{n}{clip\PYZus{}test}\PY{p}{(}\PY{n}{clip}\PY{p}{,} \PY{l+m+mi}{10}\PY{p}{)}
\PY{n}{clip\PYZus{}test}\PY{p}{(}\PY{n}{clip}\PY{p}{,} \PY{l+m+mi}{5}\PY{p}{)}
    
\end{Verbatim}
\end{tcolorbox}

    \begin{Verbatim}[commandchars=\\\{\}]

Gradients for mValue=10
gradients["dWaa"][1][2] = 10.0
gradients["dWax"][3][1] = -10.0
gradients["dWya"][1][2] = 0.2971381536101662
gradients["db"][4] = [10.]
gradients["dby"][1] = [8.45833407]
\textcolor{ansi-green-intense}{All tests passed!}

Gradients for mValue=5
gradients["dWaa"][1][2] = 5.0
gradients["dWax"][3][1] = -5.0
gradients["dWya"][1][2] = 0.2971381536101662
gradients["db"][4] = [5.]
gradients["dby"][1] = [5.]
\textcolor{ansi-green-intense}{All tests passed!}
    \end{Verbatim}

    \textbf{Expected values}

\begin{verbatim}
Gradients for mValue=10
gradients["dWaa"][1][2] = 10.0
gradients["dWax"][3][1] = -10.0
gradients["dWya"][1][2] = 0.2971381536101662
gradients["db"][4] = [10.]
gradients["dby"][1] = [8.45833407]

Gradients for mValue=5
gradients["dWaa"][1][2] = 5.0
gradients["dWax"][3][1] = -5.0
gradients["dWya"][1][2] = 0.2971381536101662
gradients["db"][4] = [5.]
gradients["dby"][1] = [5.]
\end{verbatim}

    \#\#\# 2.2 - Sampling

Now, assume that your model is trained, and you would like to generate
new text (characters). The process of generation is explained in the
picture below:

Figure 3: In this picture, you can assume the model is already trained.
You pass in \(x^{\langle 1\rangle} = \vec{0}\) at the first time-step,
and have the network sample one character at a time.

    \#\#\# Exercise 2 - sample

Implement the \texttt{sample} function below to sample characters.

You need to carry out 4 steps:

\begin{itemize}
\tightlist
\item
  \textbf{Step 1}: Input the ``dummy'' vector of zeros
  \(x^{\langle 1 \rangle} = \vec{0}\).

  \begin{itemize}
  \tightlist
  \item
    This is the default input before you've generated any characters.
    You also set \(a^{\langle 0 \rangle} = \vec{0}\)
  \end{itemize}
\end{itemize}

    \begin{itemize}
\tightlist
\item
  \textbf{Step 2}: Run one step of forward propagation to get
  \(a^{\langle 1 \rangle}\) and \(\hat{y}^{\langle 1 \rangle}\). Here
  are the equations:
\end{itemize}

\emph{hidden state:}\\
\[ a^{\langle t+1 \rangle} = \tanh(W_{ax}  x^{\langle t+1 \rangle } + W_{aa} a^{\langle t \rangle } + b)\tag{1}\]

\emph{activation:}
\[ z^{\langle t + 1 \rangle } = W_{ya}  a^{\langle t + 1 \rangle } + b_y \tag{2}\]

\emph{prediction:}
\[ \hat{y}^{\langle t+1 \rangle } = softmax(z^{\langle t + 1 \rangle })\tag{3}\]

\begin{itemize}
\tightlist
\item
  Details about \(\hat{y}^{\langle t+1 \rangle }\):

  \begin{itemize}
  \tightlist
  \item
    Note that \(\hat{y}^{\langle t+1 \rangle }\) is a (softmax)
    probability vector (its entries are between 0 and 1 and sum to 1).
  \item
    \(\hat{y}^{\langle t+1 \rangle}_i\) represents the probability that
    the character indexed by ``i'' is the next character.\\
  \item
    A \texttt{softmax()} function is provided for you to use.
  \end{itemize}
\end{itemize}

    Additional Hints: Matrix multiplication with numpy

\(x^{\langle 1 \rangle}\) is \texttt{x} in the code. When creating the
one-hot vector, make a numpy array of zeros, with the number of rows
equal to the number of unique characters, and the number of columns
equal to one. It's a 2D and not a 1D array.

\(a^{\langle 0 \rangle}\) is \texttt{a\_prev} in the code. It is a numpy
array of zeros, where the number of rows is \(n_{a}\), and number of
columns is 1. It is a 2D array as well. \(n_{a}\) is retrieved by
getting the number of columns in \(W_{aa}\) (the numbers need to match
in order for the matrix multiplication \(W_{aa}a^{\langle t \rangle}\)
to work.

Official documentation for numpy.dot and numpy.tanh.

Below is some revision code for matrix multiplication

\begin{Shaded}
\begin{Highlighting}[]
\NormalTok{matrix1 }\OperatorTok{=}\NormalTok{ np.array([[}\DecValTok{1}\NormalTok{,}\DecValTok{1}\NormalTok{],[}\DecValTok{2}\NormalTok{,}\DecValTok{2}\NormalTok{],[}\DecValTok{3}\NormalTok{,}\DecValTok{3}\NormalTok{]]) }\CommentTok{\# (3,2)}
\NormalTok{matrix2 }\OperatorTok{=}\NormalTok{ np.array([[}\DecValTok{0}\NormalTok{],[}\DecValTok{0}\NormalTok{],[}\DecValTok{0}\NormalTok{]]) }\CommentTok{\# (3,1) }
\NormalTok{vector1D }\OperatorTok{=}\NormalTok{ np.array([}\DecValTok{1}\NormalTok{,}\DecValTok{1}\NormalTok{]) }\CommentTok{\# (2,) }
\NormalTok{vector2D }\OperatorTok{=}\NormalTok{ np.array([[}\DecValTok{1}\NormalTok{],[}\DecValTok{1}\NormalTok{]]) }\CommentTok{\# (2,1)}
\BuiltInTok{print}\NormalTok{(}\StringTok{"matrix1 }\CharTok{\textbackslash{}n}\StringTok{"}\NormalTok{, matrix1,}\StringTok{"}\CharTok{\textbackslash{}n}\StringTok{"}\NormalTok{)}
\BuiltInTok{print}\NormalTok{(}\StringTok{"matrix2 }\CharTok{\textbackslash{}n}\StringTok{"}\NormalTok{, matrix2,}\StringTok{"}\CharTok{\textbackslash{}n}\StringTok{"}\NormalTok{)}
\BuiltInTok{print}\NormalTok{(}\StringTok{"vector1D }\CharTok{\textbackslash{}n}\StringTok{"}\NormalTok{, vector1D,}\StringTok{"}\CharTok{\textbackslash{}n}\StringTok{"}\NormalTok{)}
\BuiltInTok{print}\NormalTok{(}\StringTok{"vector2D }\CharTok{\textbackslash{}n}\StringTok{"}\NormalTok{, vector2D)}
\end{Highlighting}
\end{Shaded}

\begin{Shaded}
\begin{Highlighting}[]
\BuiltInTok{print}\NormalTok{(}\StringTok{"Multiply 2D and 1D arrays: result is a 1D array}\CharTok{\textbackslash{}n}\StringTok{"}\NormalTok{, }
\NormalTok{      np.dot(matrix1,vector1D))}
\BuiltInTok{print}\NormalTok{(}\StringTok{"Multiply 2D and 2D arrays: result is a 2D array}\CharTok{\textbackslash{}n}\StringTok{"}\NormalTok{, }
\NormalTok{      np.dot(matrix1,vector2D))}
\end{Highlighting}
\end{Shaded}

\begin{Shaded}
\begin{Highlighting}[]
\BuiltInTok{print}\NormalTok{(}\StringTok{"Adding (3 x 1) vector to a (3 x 1) vector is a (3 x 1) vector}\CharTok{\textbackslash{}n}\StringTok{"}\NormalTok{,}
      \StringTok{"This is what we want here!}\CharTok{\textbackslash{}n}\StringTok{"}\NormalTok{, }
\NormalTok{      np.dot(matrix1,vector2D) }\OperatorTok{+}\NormalTok{ matrix2)}
\end{Highlighting}
\end{Shaded}

\begin{Shaded}
\begin{Highlighting}[]
\BuiltInTok{print}\NormalTok{(}\StringTok{"Adding a (3,) vector to a (3 x 1) vector}\CharTok{\textbackslash{}n}\StringTok{"}\NormalTok{,}
      \StringTok{"broadcasts the 1D array across the second dimension}\CharTok{\textbackslash{}n}\StringTok{"}\NormalTok{,}
      \StringTok{"Not what we want here!}\CharTok{\textbackslash{}n}\StringTok{"}\NormalTok{,}
\NormalTok{      np.dot(matrix1,vector1D) }\OperatorTok{+}\NormalTok{ matrix2}
\NormalTok{     )}
\end{Highlighting}
\end{Shaded}

    \begin{itemize}
\item
  \textbf{Step 3}: Sampling:

  \begin{itemize}
  \tightlist
  \item
    Now that you have \(y^{\langle t+1 \rangle}\), you want to select
    the next letter in the dinosaur name. If you select the most
    probable, the model will always generate the same result given a
    starting letter. To make the results more interesting, use
    \texttt{np.random.choice} to select a next letter that is
    \emph{likely}, but not always the same.
  \item
    Pick the next character's \textbf{index} according to the
    probability distribution specified by
    \(\hat{y}^{\langle t+1 \rangle }\).
  \item
    This means that if \(\hat{y}^{\langle t+1 \rangle }_i = 0.16\), you
    will pick the index ``i'' with 16\% probability.
  \item
    Use
    \href{https://docs.scipy.org/doc/numpy-1.13.0/reference/generated/numpy.random.choice.html}{np.random.choice}.
  \end{itemize}

  Example of how to use \texttt{np.random.choice()}:
  \texttt{python\ \ \ np.random.seed(0)\ \ \ probs\ =\ np.array({[}0.1,\ 0.0,\ 0.7,\ 0.2{]})\ \ \ idx\ =\ np.random.choice(range(len(probs)),\ p\ =\ probs)}

  \begin{itemize}
  \tightlist
  \item
    This means that you will pick the index (\texttt{idx}) according to
    the distribution:
  \end{itemize}

  \(P(index = 0) = 0.1, P(index = 1) = 0.0, P(index = 2) = 0.7, P(index = 3) = 0.2\).

  \begin{itemize}
  \tightlist
  \item
    Note that the value that's set to \texttt{p} should be set to a 1D
    vector.
  \item
    Also notice that \(\hat{y}^{\langle t+1 \rangle}\), which is
    \texttt{y} in the code, is a 2D array.
  \item
    Also notice, while in your implementation, the first argument to
    \texttt{np.random.choice} is just an ordered list {[}0,1,..,
    vocab\_len-1{]}, it is \emph{not} appropriate to use
    \texttt{char\_to\_ix.values()}. The \emph{order} of values returned
    by a Python dictionary \texttt{.values()} call will be the same
    order as they are added to the dictionary. The grader may have a
    different order when it runs your routine than when you run it in
    your notebook.
  \end{itemize}
\end{itemize}

    \hypertarget{additional-hints}{%
\subparagraph{Additional Hints}\label{additional-hints}}

\begin{itemize}
\tightlist
\item
  Documentation for the built-in Python function
  \href{https://docs.python.org/3/library/functions.html\#func-range}{range}
\item
  Docs for
  \href{https://docs.scipy.org/doc/numpy/reference/generated/numpy.ravel.html}{numpy.ravel},
  which takes a multi-dimensional array and returns its contents inside
  of a 1D vector.
\end{itemize}

\begin{Shaded}
\begin{Highlighting}[]
\NormalTok{arr }\OperatorTok{=}\NormalTok{ np.array([[}\DecValTok{1}\NormalTok{,}\DecValTok{2}\NormalTok{],[}\DecValTok{3}\NormalTok{,}\DecValTok{4}\NormalTok{]])}
\BuiltInTok{print}\NormalTok{(}\StringTok{"arr"}\NormalTok{)}
\BuiltInTok{print}\NormalTok{(arr)}
\BuiltInTok{print}\NormalTok{(}\StringTok{"arr.ravel()"}\NormalTok{)}
\BuiltInTok{print}\NormalTok{(arr.ravel())}
\end{Highlighting}
\end{Shaded}

Output:

\begin{Shaded}
\begin{Highlighting}[]
\NormalTok{arr}
\NormalTok{[[}\DecValTok{1} \DecValTok{2}\NormalTok{]}
\NormalTok{ [}\DecValTok{3} \DecValTok{4}\NormalTok{]]}
\NormalTok{arr.ravel()}
\NormalTok{[}\DecValTok{1} \DecValTok{2} \DecValTok{3} \DecValTok{4}\NormalTok{]}
\end{Highlighting}
\end{Shaded}

\begin{itemize}
\tightlist
\item
  Note that \texttt{append} is an ``in-place'' operation, which means
  the changes made by the method will remain after the call completes.
  In other words, don't do this:
\end{itemize}

\begin{Shaded}
\begin{Highlighting}[]
\NormalTok{fun\_hobbies }\OperatorTok{=}\NormalTok{ fun\_hobbies.append(}\StringTok{\textquotesingle{}learning\textquotesingle{}}\NormalTok{)  }\CommentTok{\#\# Doesn\textquotesingle{}t give you what you want!}
\end{Highlighting}
\end{Shaded}

    \begin{itemize}
\tightlist
\item
  \textbf{Step 4}: Update to \(x^{\langle t \rangle }\)

  \begin{itemize}
  \tightlist
  \item
    The last step to implement in \texttt{sample()} is to update the
    variable \texttt{x}, which currently stores
    \(x^{\langle t \rangle }\), with the value of
    \(x^{\langle t + 1 \rangle }\).
  \item
    You will represent \(x^{\langle t + 1 \rangle }\) by creating a
    one-hot vector corresponding to the character that you have chosen
    as your prediction.
  \item
    You will then forward propagate \(x^{\langle t + 1 \rangle }\) in
    Step 1 and keep repeating the process until you get a
    \texttt{"\textbackslash{}n"} character, indicating that you have
    reached the end of the dinosaur name.
  \end{itemize}
\end{itemize}

    \hypertarget{additional-hints}{%
\subparagraph{Additional Hints}\label{additional-hints}}

\begin{itemize}
\tightlist
\item
  In order to reset \texttt{x} before setting it to the new one-hot
  vector, you'll want to set all the values to zero.

  \begin{itemize}
  \tightlist
  \item
    You can either create a new numpy array:
    \href{https://docs.scipy.org/doc/numpy/reference/generated/numpy.zeros.html}{numpy.zeros}
  \item
    Or fill all values with a single number:
    \href{https://docs.scipy.org/doc/numpy/reference/generated/numpy.ndarray.fill.html}{numpy.ndarray.fill}
  \end{itemize}
\end{itemize}

    \begin{tcolorbox}[breakable, size=fbox, boxrule=1pt, pad at break*=1mm,colback=cellbackground, colframe=cellborder]
\prompt{In}{incolor}{ }{\boxspacing}
\begin{Verbatim}[commandchars=\\\{\}]

\end{Verbatim}
\end{tcolorbox}

    \begin{tcolorbox}[breakable, size=fbox, boxrule=1pt, pad at break*=1mm,colback=cellbackground, colframe=cellborder]
\prompt{In}{incolor}{52}{\boxspacing}
\begin{Verbatim}[commandchars=\\\{\}]
\PY{c+c1}{\PYZsh{} UNQ\PYZus{}C2 (UNIQUE CELL IDENTIFIER, DO NOT EDIT)}
\PY{c+c1}{\PYZsh{} GRADED FUNCTION: sample}

\PY{k}{def} \PY{n+nf}{sample}\PY{p}{(}\PY{n}{parameters}\PY{p}{,} \PY{n}{char\PYZus{}to\PYZus{}ix}\PY{p}{,} \PY{n}{seed}\PY{p}{)}\PY{p}{:}
    \PY{l+s+sd}{\PYZdq{}\PYZdq{}\PYZdq{}}
\PY{l+s+sd}{    Sample a sequence of characters according to a sequence of probability distributions output of the RNN}

\PY{l+s+sd}{    Arguments:}
\PY{l+s+sd}{    parameters \PYZhy{}\PYZhy{} Python dictionary containing the parameters Waa, Wax, Wya, by, and b. }
\PY{l+s+sd}{    char\PYZus{}to\PYZus{}ix \PYZhy{}\PYZhy{} Python dictionary mapping each character to an index.}
\PY{l+s+sd}{    seed \PYZhy{}\PYZhy{} Used for grading purposes. Do not worry about it.}

\PY{l+s+sd}{    Returns:}
\PY{l+s+sd}{    indices \PYZhy{}\PYZhy{} A list of length n containing the indices of the sampled characters.}
\PY{l+s+sd}{    \PYZdq{}\PYZdq{}\PYZdq{}}
    
    \PY{c+c1}{\PYZsh{} Retrieve parameters and relevant shapes from \PYZdq{}parameters\PYZdq{} dictionary}
    \PY{n}{Waa}\PY{p}{,} \PY{n}{Wax}\PY{p}{,} \PY{n}{Wya}\PY{p}{,} \PY{n}{by}\PY{p}{,} \PY{n}{b} \PY{o}{=} \PY{n}{parameters}\PY{p}{[}\PY{l+s+s1}{\PYZsq{}}\PY{l+s+s1}{Waa}\PY{l+s+s1}{\PYZsq{}}\PY{p}{]}\PY{p}{,} \PY{n}{parameters}\PY{p}{[}\PY{l+s+s1}{\PYZsq{}}\PY{l+s+s1}{Wax}\PY{l+s+s1}{\PYZsq{}}\PY{p}{]}\PY{p}{,} \PY{n}{parameters}\PY{p}{[}\PY{l+s+s1}{\PYZsq{}}\PY{l+s+s1}{Wya}\PY{l+s+s1}{\PYZsq{}}\PY{p}{]}\PY{p}{,} \PY{n}{parameters}\PY{p}{[}\PY{l+s+s1}{\PYZsq{}}\PY{l+s+s1}{by}\PY{l+s+s1}{\PYZsq{}}\PY{p}{]}\PY{p}{,} \PY{n}{parameters}\PY{p}{[}\PY{l+s+s1}{\PYZsq{}}\PY{l+s+s1}{b}\PY{l+s+s1}{\PYZsq{}}\PY{p}{]}
    \PY{n}{vocab\PYZus{}size} \PY{o}{=} \PY{n}{by}\PY{o}{.}\PY{n}{shape}\PY{p}{[}\PY{l+m+mi}{0}\PY{p}{]}
    \PY{n}{n\PYZus{}a} \PY{o}{=} \PY{n}{Waa}\PY{o}{.}\PY{n}{shape}\PY{p}{[}\PY{l+m+mi}{1}\PY{p}{]}
    
    \PY{c+c1}{\PYZsh{}\PYZsh{}\PYZsh{} START CODE HERE \PYZsh{}\PYZsh{}\PYZsh{}}
    \PY{c+c1}{\PYZsh{} Step 1: Create the a zero vector x that can be used as the one\PYZhy{}hot vector }
    \PY{c+c1}{\PYZsh{} Representing the first character (initializing the sequence generation). (≈1 line)}
    \PY{n}{x} \PY{o}{=} \PY{n}{np}\PY{o}{.}\PY{n}{zeros}\PY{p}{(}\PY{p}{(}\PY{n}{vocab\PYZus{}size}\PY{p}{,} \PY{l+m+mi}{1}\PY{p}{)}\PY{p}{)}
    \PY{c+c1}{\PYZsh{} Step 1\PYZsq{}: Initialize a\PYZus{}prev as zeros (≈1 line)}
    \PY{n}{a\PYZus{}prev} \PY{o}{=} \PY{n}{np}\PY{o}{.}\PY{n}{zeros}\PY{p}{(}\PY{p}{(}\PY{n}{n\PYZus{}a}\PY{p}{,} \PY{l+m+mi}{1}\PY{p}{)}\PY{p}{)}
    
    \PY{c+c1}{\PYZsh{} Create an empty list of indices. This is the list which will contain the list of indices of the characters to generate (≈1 line)}
    \PY{n}{indices} \PY{o}{=} \PY{p}{[}\PY{p}{]}
    
    \PY{c+c1}{\PYZsh{} idx is the index of the one\PYZhy{}hot vector x that is set to 1}
    \PY{c+c1}{\PYZsh{} All other positions in x are zero.}
    \PY{c+c1}{\PYZsh{} Initialize idx to \PYZhy{}1}
    \PY{n}{idx} \PY{o}{=} \PY{o}{\PYZhy{}}\PY{l+m+mi}{1} 
    
    \PY{c+c1}{\PYZsh{} Loop over time\PYZhy{}steps t. At each time\PYZhy{}step:}
    \PY{c+c1}{\PYZsh{} Sample a character from a probability distribution }
    \PY{c+c1}{\PYZsh{} And append its index (`idx`) to the list \PYZdq{}indices\PYZdq{}. }
    \PY{c+c1}{\PYZsh{} You\PYZsq{}ll stop if you reach 50 characters }
    \PY{c+c1}{\PYZsh{} (which should be very unlikely with a well\PYZhy{}trained model).}
    \PY{c+c1}{\PYZsh{} Setting the maximum number of characters helps with debugging and prevents infinite loops. }
    \PY{n}{counter} \PY{o}{=} \PY{l+m+mi}{0}
    \PY{n}{newline\PYZus{}character} \PY{o}{=} \PY{n}{char\PYZus{}to\PYZus{}ix}\PY{p}{[}\PY{l+s+s1}{\PYZsq{}}\PY{l+s+se}{\PYZbs{}n}\PY{l+s+s1}{\PYZsq{}}\PY{p}{]}
    
    \PY{k}{while} \PY{p}{(}\PY{n}{idx} \PY{o}{!=} \PY{n}{newline\PYZus{}character} \PY{o+ow}{and} \PY{n}{counter} \PY{o}{!=} \PY{l+m+mi}{50}\PY{p}{)}\PY{p}{:}
        
        \PY{c+c1}{\PYZsh{} Step 2: Forward propagate x using the equations (1), (2) and (3)}
        \PY{n}{a} \PY{o}{=} \PY{n}{np}\PY{o}{.}\PY{n}{tanh}\PY{p}{(}\PY{n}{np}\PY{o}{.}\PY{n}{matmul}\PY{p}{(}\PY{n}{Wax}\PY{p}{,} \PY{n}{x}\PY{p}{)} \PY{o}{+} \PY{n}{np}\PY{o}{.}\PY{n}{matmul}\PY{p}{(}\PY{n}{Waa}\PY{p}{,} \PY{n}{a\PYZus{}prev}\PY{p}{)} \PY{o}{+} \PY{n}{b}\PY{p}{)}
        \PY{n}{z} \PY{o}{=} \PY{n}{np}\PY{o}{.}\PY{n}{matmul}\PY{p}{(}\PY{n}{Wya}\PY{p}{,} \PY{n}{a}\PY{p}{)} \PY{o}{+} \PY{n}{by}
        \PY{n}{y} \PY{o}{=} \PY{n}{softmax}\PY{p}{(}\PY{n}{z}\PY{p}{)}
        
        \PY{c+c1}{\PYZsh{} For grading purposes}
        \PY{n}{np}\PY{o}{.}\PY{n}{random}\PY{o}{.}\PY{n}{seed}\PY{p}{(}\PY{n}{counter} \PY{o}{+} \PY{n}{seed}\PY{p}{)} 
        
        \PY{c+c1}{\PYZsh{} Step 3: Sample the index of a character within the vocabulary from the probability distribution y}
        \PY{c+c1}{\PYZsh{} (see additional hints above)}
        \PY{n}{idx} \PY{o}{=} \PY{n}{np}\PY{o}{.}\PY{n}{random}\PY{o}{.}\PY{n}{choice}\PY{p}{(}\PY{n+nb}{range}\PY{p}{(}\PY{n}{vocab\PYZus{}size}\PY{p}{)}\PY{p}{,} \PY{n}{p} \PY{o}{=} \PY{n}{y}\PY{p}{[}\PY{p}{:}\PY{p}{,}\PY{l+m+mi}{0}\PY{p}{]}\PY{p}{)}

        \PY{c+c1}{\PYZsh{} Append the index to \PYZdq{}indices\PYZdq{}}
        \PY{n}{indices}\PY{o}{.}\PY{n}{append}\PY{p}{(}\PY{n}{idx}\PY{p}{)}
        
        \PY{c+c1}{\PYZsh{} Step 4: Overwrite the input x with one that corresponds to the sampled index `idx`.}
        \PY{c+c1}{\PYZsh{} (see additional hints above)}
        \PY{n}{x} \PY{o}{=} \PY{n}{np}\PY{o}{.}\PY{n}{zeros}\PY{p}{(}\PY{p}{(}\PY{n}{vocab\PYZus{}size}\PY{p}{,} \PY{l+m+mi}{1}\PY{p}{)}\PY{p}{)}
        \PY{n}{x}\PY{p}{[}\PY{n}{idx}\PY{p}{]} \PY{o}{=} \PY{l+m+mi}{1}
        
        \PY{c+c1}{\PYZsh{} Update \PYZdq{}a\PYZus{}prev\PYZdq{} to be \PYZdq{}a\PYZdq{}}
        \PY{n}{a\PYZus{}prev} \PY{o}{=} \PY{n}{a}
        
        \PY{c+c1}{\PYZsh{} for grading purposes}
        \PY{n}{seed} \PY{o}{+}\PY{o}{=} \PY{l+m+mi}{1}
        \PY{n}{counter} \PY{o}{+}\PY{o}{=}\PY{l+m+mi}{1}
        
    \PY{c+c1}{\PYZsh{}\PYZsh{}\PYZsh{} END CODE HERE \PYZsh{}\PYZsh{}\PYZsh{}}

    \PY{k}{if} \PY{p}{(}\PY{n}{counter} \PY{o}{==} \PY{l+m+mi}{50}\PY{p}{)}\PY{p}{:}
        \PY{n}{indices}\PY{o}{.}\PY{n}{append}\PY{p}{(}\PY{n}{char\PYZus{}to\PYZus{}ix}\PY{p}{[}\PY{l+s+s1}{\PYZsq{}}\PY{l+s+se}{\PYZbs{}n}\PY{l+s+s1}{\PYZsq{}}\PY{p}{]}\PY{p}{)}
    
    \PY{k}{return} \PY{n}{indices}
\end{Verbatim}
\end{tcolorbox}

    \begin{tcolorbox}[breakable, size=fbox, boxrule=1pt, pad at break*=1mm,colback=cellbackground, colframe=cellborder]
\prompt{In}{incolor}{53}{\boxspacing}
\begin{Verbatim}[commandchars=\\\{\}]
\PY{k}{def} \PY{n+nf}{sample\PYZus{}test}\PY{p}{(}\PY{n}{target}\PY{p}{)}\PY{p}{:}
    \PY{n}{np}\PY{o}{.}\PY{n}{random}\PY{o}{.}\PY{n}{seed}\PY{p}{(}\PY{l+m+mi}{24}\PY{p}{)}
    \PY{n}{\PYZus{}}\PY{p}{,} \PY{n}{n\PYZus{}a} \PY{o}{=} \PY{l+m+mi}{20}\PY{p}{,} \PY{l+m+mi}{100}
    \PY{n}{Wax}\PY{p}{,} \PY{n}{Waa}\PY{p}{,} \PY{n}{Wya} \PY{o}{=} \PY{n}{np}\PY{o}{.}\PY{n}{random}\PY{o}{.}\PY{n}{randn}\PY{p}{(}\PY{n}{n\PYZus{}a}\PY{p}{,} \PY{n}{vocab\PYZus{}size}\PY{p}{)}\PY{p}{,} \PY{n}{np}\PY{o}{.}\PY{n}{random}\PY{o}{.}\PY{n}{randn}\PY{p}{(}\PY{n}{n\PYZus{}a}\PY{p}{,} \PY{n}{n\PYZus{}a}\PY{p}{)}\PY{p}{,} \PY{n}{np}\PY{o}{.}\PY{n}{random}\PY{o}{.}\PY{n}{randn}\PY{p}{(}\PY{n}{vocab\PYZus{}size}\PY{p}{,} \PY{n}{n\PYZus{}a}\PY{p}{)}
    \PY{n}{b}\PY{p}{,} \PY{n}{by} \PY{o}{=} \PY{n}{np}\PY{o}{.}\PY{n}{random}\PY{o}{.}\PY{n}{randn}\PY{p}{(}\PY{n}{n\PYZus{}a}\PY{p}{,} \PY{l+m+mi}{1}\PY{p}{)}\PY{p}{,} \PY{n}{np}\PY{o}{.}\PY{n}{random}\PY{o}{.}\PY{n}{randn}\PY{p}{(}\PY{n}{vocab\PYZus{}size}\PY{p}{,} \PY{l+m+mi}{1}\PY{p}{)}
    \PY{n}{parameters} \PY{o}{=} \PY{p}{\PYZob{}}\PY{l+s+s2}{\PYZdq{}}\PY{l+s+s2}{Wax}\PY{l+s+s2}{\PYZdq{}}\PY{p}{:} \PY{n}{Wax}\PY{p}{,} \PY{l+s+s2}{\PYZdq{}}\PY{l+s+s2}{Waa}\PY{l+s+s2}{\PYZdq{}}\PY{p}{:} \PY{n}{Waa}\PY{p}{,} \PY{l+s+s2}{\PYZdq{}}\PY{l+s+s2}{Wya}\PY{l+s+s2}{\PYZdq{}}\PY{p}{:} \PY{n}{Wya}\PY{p}{,} \PY{l+s+s2}{\PYZdq{}}\PY{l+s+s2}{b}\PY{l+s+s2}{\PYZdq{}}\PY{p}{:} \PY{n}{b}\PY{p}{,} \PY{l+s+s2}{\PYZdq{}}\PY{l+s+s2}{by}\PY{l+s+s2}{\PYZdq{}}\PY{p}{:} \PY{n}{by}\PY{p}{\PYZcb{}}


    \PY{n}{indices} \PY{o}{=} \PY{n}{target}\PY{p}{(}\PY{n}{parameters}\PY{p}{,} \PY{n}{char\PYZus{}to\PYZus{}ix}\PY{p}{,} \PY{l+m+mi}{0}\PY{p}{)}
    \PY{n+nb}{print}\PY{p}{(}\PY{l+s+s2}{\PYZdq{}}\PY{l+s+s2}{Sampling:}\PY{l+s+s2}{\PYZdq{}}\PY{p}{)}
    \PY{n+nb}{print}\PY{p}{(}\PY{l+s+s2}{\PYZdq{}}\PY{l+s+s2}{list of sampled indices:}\PY{l+s+se}{\PYZbs{}n}\PY{l+s+s2}{\PYZdq{}}\PY{p}{,} \PY{n}{indices}\PY{p}{)}
    \PY{n+nb}{print}\PY{p}{(}\PY{l+s+s2}{\PYZdq{}}\PY{l+s+s2}{list of sampled characters:}\PY{l+s+se}{\PYZbs{}n}\PY{l+s+s2}{\PYZdq{}}\PY{p}{,} \PY{p}{[}\PY{n}{ix\PYZus{}to\PYZus{}char}\PY{p}{[}\PY{n}{i}\PY{p}{]} \PY{k}{for} \PY{n}{i} \PY{o+ow}{in} \PY{n}{indices}\PY{p}{]}\PY{p}{)}
    
    \PY{k}{assert} \PY{n+nb}{len}\PY{p}{(}\PY{n}{indices}\PY{p}{)} \PY{o}{\PYZlt{}} \PY{l+m+mi}{52}\PY{p}{,} \PY{l+s+s2}{\PYZdq{}}\PY{l+s+s2}{Indices length must be smaller than 52}\PY{l+s+s2}{\PYZdq{}}
    \PY{k}{assert} \PY{n}{indices}\PY{p}{[}\PY{o}{\PYZhy{}}\PY{l+m+mi}{1}\PY{p}{]} \PY{o}{==} \PY{n}{char\PYZus{}to\PYZus{}ix}\PY{p}{[}\PY{l+s+s1}{\PYZsq{}}\PY{l+s+se}{\PYZbs{}n}\PY{l+s+s1}{\PYZsq{}}\PY{p}{]}\PY{p}{,} \PY{l+s+s2}{\PYZdq{}}\PY{l+s+s2}{All samples must end with }\PY{l+s+se}{\PYZbs{}\PYZbs{}}\PY{l+s+s2}{n}\PY{l+s+s2}{\PYZdq{}}
    \PY{k}{assert} \PY{n+nb}{min}\PY{p}{(}\PY{n}{indices}\PY{p}{)} \PY{o}{\PYZgt{}}\PY{o}{=} \PY{l+m+mi}{0} \PY{o+ow}{and} \PY{n+nb}{max}\PY{p}{(}\PY{n}{indices}\PY{p}{)} \PY{o}{\PYZlt{}} \PY{n+nb}{len}\PY{p}{(}\PY{n}{char\PYZus{}to\PYZus{}ix}\PY{p}{)}\PY{p}{,} \PY{l+s+sa}{f}\PY{l+s+s2}{\PYZdq{}}\PY{l+s+s2}{Sampled indexes must be between 0 and len(char\PYZus{}to\PYZus{}ix)=}\PY{l+s+si}{\PYZob{}}\PY{n+nb}{len}\PY{p}{(}\PY{n}{char\PYZus{}to\PYZus{}ix}\PY{p}{)}\PY{l+s+si}{\PYZcb{}}\PY{l+s+s2}{\PYZdq{}}
    \PY{k}{assert} \PY{n}{np}\PY{o}{.}\PY{n}{allclose}\PY{p}{(}\PY{n}{indices}\PY{p}{[}\PY{l+m+mi}{0}\PY{p}{:}\PY{l+m+mi}{6}\PY{p}{]}\PY{p}{,} \PY{p}{[}\PY{l+m+mi}{23}\PY{p}{,} \PY{l+m+mi}{16}\PY{p}{,} \PY{l+m+mi}{26}\PY{p}{,} \PY{l+m+mi}{26}\PY{p}{,} \PY{l+m+mi}{24}\PY{p}{,} \PY{l+m+mi}{3}\PY{p}{]}\PY{p}{)}\PY{p}{,} \PY{l+s+s2}{\PYZdq{}}\PY{l+s+s2}{Wrong values}\PY{l+s+s2}{\PYZdq{}}
    
    \PY{n+nb}{print}\PY{p}{(}\PY{l+s+s2}{\PYZdq{}}\PY{l+s+se}{\PYZbs{}033}\PY{l+s+s2}{[92mAll tests passed!}\PY{l+s+s2}{\PYZdq{}}\PY{p}{)}

\PY{n}{sample\PYZus{}test}\PY{p}{(}\PY{n}{sample}\PY{p}{)}
\end{Verbatim}
\end{tcolorbox}

    \begin{Verbatim}[commandchars=\\\{\}]
Sampling:
list of sampled indices:
 [23, 16, 26, 26, 24, 3, 21, 1, 7, 24, 15, 3, 25, 20, 6, 13, 10, 8, 20, 12, 2,
0]
list of sampled characters:
 ['w', 'p', 'z', 'z', 'x', 'c', 'u', 'a', 'g', 'x', 'o', 'c', 'y', 't', 'f',
'm', 'j', 'h', 't', 'l', 'b', '\textbackslash{}n']
\textcolor{ansi-green-intense}{All tests passed!}
    \end{Verbatim}

    \textbf{Expected output}

\begin{verbatim}
Sampling:
list of sampled indices:
 [23, 16, 26, 26, 24, 3, 21, 1, 7, 24, 15, 3, 25, 20, 6, 13, 10, 8, 20, 12, 2, 0]
list of sampled characters:
 ['w', 'p', 'z', 'z', 'x', 'c', 'u', 'a', 'g', 'x', 'o', 'c', 'y', 't', 'f', 'm', 'j', 'h', 't', 'l', 'b', '\n']
\end{verbatim}

    What you should remember: * Very large, or ``exploding'' gradients
updates can be so large that they ``overshoot'' the optimal values
during back prop -- making training difficult * Clip gradients before
updating the parameters to avoid exploding gradients * Sampling is a
technique you can use to pick the index of the next character according
to a probability distribution. * To begin character-level sampling: *
Input a ``dummy'' vector of zeros as a default input * Run one step of
forward propagation to get 𝑎⟨1⟩ (your first character) and 𝑦̂ ⟨1⟩
(probability distribution for the following character) * When sampling,
avoid generating the same result each time given the starting letter
(and make your names more interesting!) by using
\texttt{np.random.choice}

    \#\# 3 - Building the Language Model

It's time to build the character-level language model for text
generation!

\#\#\# 3.1 - Gradient Descent

In this section you will implement a function performing one step of
stochastic gradient descent (with clipped gradients). You'll go through
the training examples one at a time, so the optimization algorithm will
be stochastic gradient descent.

As a reminder, here are the steps of a common optimization loop for an
RNN:

\begin{itemize}
\tightlist
\item
  Forward propagate through the RNN to compute the loss
\item
  Backward propagate through time to compute the gradients of the loss
  with respect to the parameters
\item
  Clip the gradients
\item
  Update the parameters using gradient descent
\end{itemize}

\#\#\# Exercise 3 - optimize

Implement the optimization process (one step of stochastic gradient
descent).

The following functions are provided:

\begin{Shaded}
\begin{Highlighting}[]
\KeywordTok{def}\NormalTok{ rnn\_forward(X, Y, a\_prev, parameters):}
    \CommentTok{""" Performs the forward propagation through the RNN and computes the cross{-}entropy loss.}
\CommentTok{    It returns the loss\textquotesingle{} value as well as a "cache" storing values to be used in backpropagation."""}
\NormalTok{    ....}
    \ControlFlowTok{return}\NormalTok{ loss, cache}
    
\KeywordTok{def}\NormalTok{ rnn\_backward(X, Y, parameters, cache):}
    \CommentTok{""" Performs the backward propagation through time to compute the gradients of the loss with respect}
\CommentTok{    to the parameters. It returns also all the hidden states."""}
\NormalTok{    ...}
    \ControlFlowTok{return}\NormalTok{ gradients, a}

\KeywordTok{def}\NormalTok{ update\_parameters(parameters, gradients, learning\_rate):}
    \CommentTok{""" Updates parameters using the Gradient Descent Update Rule."""}
\NormalTok{    ...}
    \ControlFlowTok{return}\NormalTok{ parameters}
\end{Highlighting}
\end{Shaded}

Recall that you previously implemented the \texttt{clip} function:

    \hypertarget{parameters}{%
\paragraph{Parameters}\label{parameters}}

\begin{itemize}
\tightlist
\item
  Note that the weights and biases inside the \texttt{parameters}
  dictionary are being updated by the optimization, even though
  \texttt{parameters} is not one of the returned values of the
  \texttt{optimize} function. The \texttt{parameters} dictionary is
  passed by reference into the function, so changes to this dictionary
  are making changes to the \texttt{parameters} dictionary even when
  accessed outside of the function.
\item
  Python dictionaries and lists are ``pass by reference'', which means
  that if you pass a dictionary into a function and modify the
  dictionary within the function, this changes that same dictionary
  (it's not a copy of the dictionary).
\end{itemize}

    \begin{tcolorbox}[breakable, size=fbox, boxrule=1pt, pad at break*=1mm,colback=cellbackground, colframe=cellborder]
\prompt{In}{incolor}{56}{\boxspacing}
\begin{Verbatim}[commandchars=\\\{\}]
\PY{c+c1}{\PYZsh{} UNQ\PYZus{}C3 (UNIQUE CELL IDENTIFIER, DO NOT EDIT)}
\PY{c+c1}{\PYZsh{} GRADED FUNCTION: optimize}

\PY{k}{def} \PY{n+nf}{optimize}\PY{p}{(}\PY{n}{X}\PY{p}{,} \PY{n}{Y}\PY{p}{,} \PY{n}{a\PYZus{}prev}\PY{p}{,} \PY{n}{parameters}\PY{p}{,} \PY{n}{learning\PYZus{}rate} \PY{o}{=} \PY{l+m+mf}{0.01}\PY{p}{)}\PY{p}{:}
    \PY{l+s+sd}{\PYZdq{}\PYZdq{}\PYZdq{}}
\PY{l+s+sd}{    Execute one step of the optimization to train the model.}
\PY{l+s+sd}{    }
\PY{l+s+sd}{    Arguments:}
\PY{l+s+sd}{    X \PYZhy{}\PYZhy{} list of integers, where each integer is a number that maps to a character in the vocabulary.}
\PY{l+s+sd}{    Y \PYZhy{}\PYZhy{} list of integers, exactly the same as X but shifted one index to the left.}
\PY{l+s+sd}{    a\PYZus{}prev \PYZhy{}\PYZhy{} previous hidden state.}
\PY{l+s+sd}{    parameters \PYZhy{}\PYZhy{} python dictionary containing:}
\PY{l+s+sd}{                        Wax \PYZhy{}\PYZhy{} Weight matrix multiplying the input, numpy array of shape (n\PYZus{}a, n\PYZus{}x)}
\PY{l+s+sd}{                        Waa \PYZhy{}\PYZhy{} Weight matrix multiplying the hidden state, numpy array of shape (n\PYZus{}a, n\PYZus{}a)}
\PY{l+s+sd}{                        Wya \PYZhy{}\PYZhy{} Weight matrix relating the hidden\PYZhy{}state to the output, numpy array of shape (n\PYZus{}y, n\PYZus{}a)}
\PY{l+s+sd}{                        b \PYZhy{}\PYZhy{}  Bias, numpy array of shape (n\PYZus{}a, 1)}
\PY{l+s+sd}{                        by \PYZhy{}\PYZhy{} Bias relating the hidden\PYZhy{}state to the output, numpy array of shape (n\PYZus{}y, 1)}
\PY{l+s+sd}{    learning\PYZus{}rate \PYZhy{}\PYZhy{} learning rate for the model.}
\PY{l+s+sd}{    }
\PY{l+s+sd}{    Returns:}
\PY{l+s+sd}{    loss \PYZhy{}\PYZhy{} value of the loss function (cross\PYZhy{}entropy)}
\PY{l+s+sd}{    gradients \PYZhy{}\PYZhy{} python dictionary containing:}
\PY{l+s+sd}{                        dWax \PYZhy{}\PYZhy{} Gradients of input\PYZhy{}to\PYZhy{}hidden weights, of shape (n\PYZus{}a, n\PYZus{}x)}
\PY{l+s+sd}{                        dWaa \PYZhy{}\PYZhy{} Gradients of hidden\PYZhy{}to\PYZhy{}hidden weights, of shape (n\PYZus{}a, n\PYZus{}a)}
\PY{l+s+sd}{                        dWya \PYZhy{}\PYZhy{} Gradients of hidden\PYZhy{}to\PYZhy{}output weights, of shape (n\PYZus{}y, n\PYZus{}a)}
\PY{l+s+sd}{                        db \PYZhy{}\PYZhy{} Gradients of bias vector, of shape (n\PYZus{}a, 1)}
\PY{l+s+sd}{                        dby \PYZhy{}\PYZhy{} Gradients of output bias vector, of shape (n\PYZus{}y, 1)}
\PY{l+s+sd}{    a[len(X)\PYZhy{}1] \PYZhy{}\PYZhy{} the last hidden state, of shape (n\PYZus{}a, 1)}
\PY{l+s+sd}{    \PYZdq{}\PYZdq{}\PYZdq{}}
    
    \PY{c+c1}{\PYZsh{}\PYZsh{}\PYZsh{} START CODE HERE \PYZsh{}\PYZsh{}\PYZsh{}}
    
    \PY{c+c1}{\PYZsh{} Forward propagate through time (≈1 line)}
    \PY{n}{loss}\PY{p}{,} \PY{n}{cache} \PY{o}{=} \PY{n}{rnn\PYZus{}forward}\PY{p}{(}\PY{n}{X}\PY{p}{,} \PY{n}{Y}\PY{p}{,} \PY{n}{a\PYZus{}prev}\PY{p}{,} \PY{n}{parameters}\PY{p}{)}
    
    \PY{c+c1}{\PYZsh{} Backpropagate through time (≈1 line)}
    \PY{n}{gradients}\PY{p}{,} \PY{n}{a} \PY{o}{=} \PY{n}{rnn\PYZus{}backward}\PY{p}{(}\PY{n}{X}\PY{p}{,} \PY{n}{Y}\PY{p}{,} \PY{n}{parameters}\PY{p}{,} \PY{n}{cache}\PY{p}{)}
    
    \PY{c+c1}{\PYZsh{} Clip your gradients between \PYZhy{}5 (min) and 5 (max) (≈1 line)}
    \PY{n}{gradients} \PY{o}{=} \PY{n}{clip}\PY{p}{(}\PY{n}{gradients}\PY{p}{,} \PY{l+m+mi}{5}\PY{p}{)}
    
    \PY{c+c1}{\PYZsh{} Update parameters (≈1 line)}
    \PY{n}{parameters} \PY{o}{=} \PY{n}{update\PYZus{}parameters}\PY{p}{(}\PY{n}{parameters}\PY{p}{,} \PY{n}{gradients}\PY{p}{,} \PY{n}{learning\PYZus{}rate}\PY{p}{)}
    
    \PY{c+c1}{\PYZsh{}\PYZsh{}\PYZsh{} END CODE HERE \PYZsh{}\PYZsh{}\PYZsh{}}
    
    \PY{k}{return} \PY{n}{loss}\PY{p}{,} \PY{n}{gradients}\PY{p}{,} \PY{n}{a}\PY{p}{[}\PY{n+nb}{len}\PY{p}{(}\PY{n}{X}\PY{p}{)}\PY{o}{\PYZhy{}}\PY{l+m+mi}{1}\PY{p}{]}
\end{Verbatim}
\end{tcolorbox}

    \begin{tcolorbox}[breakable, size=fbox, boxrule=1pt, pad at break*=1mm,colback=cellbackground, colframe=cellborder]
\prompt{In}{incolor}{57}{\boxspacing}
\begin{Verbatim}[commandchars=\\\{\}]
\PY{k}{def} \PY{n+nf}{optimize\PYZus{}test}\PY{p}{(}\PY{n}{target}\PY{p}{)}\PY{p}{:}
    \PY{n}{np}\PY{o}{.}\PY{n}{random}\PY{o}{.}\PY{n}{seed}\PY{p}{(}\PY{l+m+mi}{1}\PY{p}{)}
    \PY{n}{vocab\PYZus{}size}\PY{p}{,} \PY{n}{n\PYZus{}a} \PY{o}{=} \PY{l+m+mi}{27}\PY{p}{,} \PY{l+m+mi}{100}
    \PY{n}{a\PYZus{}prev} \PY{o}{=} \PY{n}{np}\PY{o}{.}\PY{n}{random}\PY{o}{.}\PY{n}{randn}\PY{p}{(}\PY{n}{n\PYZus{}a}\PY{p}{,} \PY{l+m+mi}{1}\PY{p}{)}
    \PY{n}{Wax}\PY{p}{,} \PY{n}{Waa}\PY{p}{,} \PY{n}{Wya} \PY{o}{=} \PY{n}{np}\PY{o}{.}\PY{n}{random}\PY{o}{.}\PY{n}{randn}\PY{p}{(}\PY{n}{n\PYZus{}a}\PY{p}{,} \PY{n}{vocab\PYZus{}size}\PY{p}{)}\PY{p}{,} \PY{n}{np}\PY{o}{.}\PY{n}{random}\PY{o}{.}\PY{n}{randn}\PY{p}{(}\PY{n}{n\PYZus{}a}\PY{p}{,} \PY{n}{n\PYZus{}a}\PY{p}{)}\PY{p}{,} \PY{n}{np}\PY{o}{.}\PY{n}{random}\PY{o}{.}\PY{n}{randn}\PY{p}{(}\PY{n}{vocab\PYZus{}size}\PY{p}{,} \PY{n}{n\PYZus{}a}\PY{p}{)}
    \PY{n}{b}\PY{p}{,} \PY{n}{by} \PY{o}{=} \PY{n}{np}\PY{o}{.}\PY{n}{random}\PY{o}{.}\PY{n}{randn}\PY{p}{(}\PY{n}{n\PYZus{}a}\PY{p}{,} \PY{l+m+mi}{1}\PY{p}{)}\PY{p}{,} \PY{n}{np}\PY{o}{.}\PY{n}{random}\PY{o}{.}\PY{n}{randn}\PY{p}{(}\PY{n}{vocab\PYZus{}size}\PY{p}{,} \PY{l+m+mi}{1}\PY{p}{)}
    \PY{n}{parameters} \PY{o}{=} \PY{p}{\PYZob{}}\PY{l+s+s2}{\PYZdq{}}\PY{l+s+s2}{Wax}\PY{l+s+s2}{\PYZdq{}}\PY{p}{:} \PY{n}{Wax}\PY{p}{,} \PY{l+s+s2}{\PYZdq{}}\PY{l+s+s2}{Waa}\PY{l+s+s2}{\PYZdq{}}\PY{p}{:} \PY{n}{Waa}\PY{p}{,} \PY{l+s+s2}{\PYZdq{}}\PY{l+s+s2}{Wya}\PY{l+s+s2}{\PYZdq{}}\PY{p}{:} \PY{n}{Wya}\PY{p}{,} \PY{l+s+s2}{\PYZdq{}}\PY{l+s+s2}{b}\PY{l+s+s2}{\PYZdq{}}\PY{p}{:} \PY{n}{b}\PY{p}{,} \PY{l+s+s2}{\PYZdq{}}\PY{l+s+s2}{by}\PY{l+s+s2}{\PYZdq{}}\PY{p}{:} \PY{n}{by}\PY{p}{\PYZcb{}}
    \PY{n}{X} \PY{o}{=} \PY{p}{[}\PY{l+m+mi}{12}\PY{p}{,} \PY{l+m+mi}{3}\PY{p}{,} \PY{l+m+mi}{5}\PY{p}{,} \PY{l+m+mi}{11}\PY{p}{,} \PY{l+m+mi}{22}\PY{p}{,} \PY{l+m+mi}{3}\PY{p}{]}
    \PY{n}{Y} \PY{o}{=} \PY{p}{[}\PY{l+m+mi}{4}\PY{p}{,} \PY{l+m+mi}{14}\PY{p}{,} \PY{l+m+mi}{11}\PY{p}{,} \PY{l+m+mi}{22}\PY{p}{,} \PY{l+m+mi}{25}\PY{p}{,} \PY{l+m+mi}{26}\PY{p}{]}
    \PY{n}{old\PYZus{}parameters} \PY{o}{=} \PY{n}{copy}\PY{o}{.}\PY{n}{deepcopy}\PY{p}{(}\PY{n}{parameters}\PY{p}{)}
    \PY{n}{loss}\PY{p}{,} \PY{n}{gradients}\PY{p}{,} \PY{n}{a\PYZus{}last} \PY{o}{=} \PY{n}{target}\PY{p}{(}\PY{n}{X}\PY{p}{,} \PY{n}{Y}\PY{p}{,} \PY{n}{a\PYZus{}prev}\PY{p}{,} \PY{n}{parameters}\PY{p}{,} \PY{n}{learning\PYZus{}rate} \PY{o}{=} \PY{l+m+mf}{0.01}\PY{p}{)}
    \PY{n+nb}{print}\PY{p}{(}\PY{l+s+s2}{\PYZdq{}}\PY{l+s+s2}{Loss =}\PY{l+s+s2}{\PYZdq{}}\PY{p}{,} \PY{n}{loss}\PY{p}{)}
    \PY{n+nb}{print}\PY{p}{(}\PY{l+s+s2}{\PYZdq{}}\PY{l+s+s2}{gradients[}\PY{l+s+se}{\PYZbs{}\PYZdq{}}\PY{l+s+s2}{dWaa}\PY{l+s+se}{\PYZbs{}\PYZdq{}}\PY{l+s+s2}{][1][2] =}\PY{l+s+s2}{\PYZdq{}}\PY{p}{,} \PY{n}{gradients}\PY{p}{[}\PY{l+s+s2}{\PYZdq{}}\PY{l+s+s2}{dWaa}\PY{l+s+s2}{\PYZdq{}}\PY{p}{]}\PY{p}{[}\PY{l+m+mi}{1}\PY{p}{]}\PY{p}{[}\PY{l+m+mi}{2}\PY{p}{]}\PY{p}{)}
    \PY{n+nb}{print}\PY{p}{(}\PY{l+s+s2}{\PYZdq{}}\PY{l+s+s2}{np.argmax(gradients[}\PY{l+s+se}{\PYZbs{}\PYZdq{}}\PY{l+s+s2}{dWax}\PY{l+s+se}{\PYZbs{}\PYZdq{}}\PY{l+s+s2}{]) =}\PY{l+s+s2}{\PYZdq{}}\PY{p}{,} \PY{n}{np}\PY{o}{.}\PY{n}{argmax}\PY{p}{(}\PY{n}{gradients}\PY{p}{[}\PY{l+s+s2}{\PYZdq{}}\PY{l+s+s2}{dWax}\PY{l+s+s2}{\PYZdq{}}\PY{p}{]}\PY{p}{)}\PY{p}{)}
    \PY{n+nb}{print}\PY{p}{(}\PY{l+s+s2}{\PYZdq{}}\PY{l+s+s2}{gradients[}\PY{l+s+se}{\PYZbs{}\PYZdq{}}\PY{l+s+s2}{dWya}\PY{l+s+se}{\PYZbs{}\PYZdq{}}\PY{l+s+s2}{][1][2] =}\PY{l+s+s2}{\PYZdq{}}\PY{p}{,} \PY{n}{gradients}\PY{p}{[}\PY{l+s+s2}{\PYZdq{}}\PY{l+s+s2}{dWya}\PY{l+s+s2}{\PYZdq{}}\PY{p}{]}\PY{p}{[}\PY{l+m+mi}{1}\PY{p}{]}\PY{p}{[}\PY{l+m+mi}{2}\PY{p}{]}\PY{p}{)}
    \PY{n+nb}{print}\PY{p}{(}\PY{l+s+s2}{\PYZdq{}}\PY{l+s+s2}{gradients[}\PY{l+s+se}{\PYZbs{}\PYZdq{}}\PY{l+s+s2}{db}\PY{l+s+se}{\PYZbs{}\PYZdq{}}\PY{l+s+s2}{][4] =}\PY{l+s+s2}{\PYZdq{}}\PY{p}{,} \PY{n}{gradients}\PY{p}{[}\PY{l+s+s2}{\PYZdq{}}\PY{l+s+s2}{db}\PY{l+s+s2}{\PYZdq{}}\PY{p}{]}\PY{p}{[}\PY{l+m+mi}{4}\PY{p}{]}\PY{p}{)}
    \PY{n+nb}{print}\PY{p}{(}\PY{l+s+s2}{\PYZdq{}}\PY{l+s+s2}{gradients[}\PY{l+s+se}{\PYZbs{}\PYZdq{}}\PY{l+s+s2}{dby}\PY{l+s+se}{\PYZbs{}\PYZdq{}}\PY{l+s+s2}{][1] =}\PY{l+s+s2}{\PYZdq{}}\PY{p}{,} \PY{n}{gradients}\PY{p}{[}\PY{l+s+s2}{\PYZdq{}}\PY{l+s+s2}{dby}\PY{l+s+s2}{\PYZdq{}}\PY{p}{]}\PY{p}{[}\PY{l+m+mi}{1}\PY{p}{]}\PY{p}{)}
    \PY{n+nb}{print}\PY{p}{(}\PY{l+s+s2}{\PYZdq{}}\PY{l+s+s2}{a\PYZus{}last[4] =}\PY{l+s+s2}{\PYZdq{}}\PY{p}{,} \PY{n}{a\PYZus{}last}\PY{p}{[}\PY{l+m+mi}{4}\PY{p}{]}\PY{p}{)}
    
    \PY{k}{assert} \PY{n}{np}\PY{o}{.}\PY{n}{isclose}\PY{p}{(}\PY{n}{loss}\PY{p}{,} \PY{l+m+mf}{126.5039757}\PY{p}{)}\PY{p}{,} \PY{l+s+s2}{\PYZdq{}}\PY{l+s+s2}{Problems with the call of the rnn\PYZus{}forward function}\PY{l+s+s2}{\PYZdq{}}
    \PY{k}{for} \PY{n}{grad} \PY{o+ow}{in} \PY{n}{gradients}\PY{o}{.}\PY{n}{values}\PY{p}{(}\PY{p}{)}\PY{p}{:}
        \PY{k}{assert} \PY{n}{np}\PY{o}{.}\PY{n}{min}\PY{p}{(}\PY{n}{grad}\PY{p}{)} \PY{o}{\PYZgt{}}\PY{o}{=} \PY{o}{\PYZhy{}}\PY{l+m+mi}{5}\PY{p}{,} \PY{l+s+s2}{\PYZdq{}}\PY{l+s+s2}{Problems in the clip function call}\PY{l+s+s2}{\PYZdq{}}
        \PY{k}{assert} \PY{n}{np}\PY{o}{.}\PY{n}{max}\PY{p}{(}\PY{n}{grad}\PY{p}{)} \PY{o}{\PYZlt{}}\PY{o}{=} \PY{l+m+mi}{5}\PY{p}{,} \PY{l+s+s2}{\PYZdq{}}\PY{l+s+s2}{Problems in the clip function call}\PY{l+s+s2}{\PYZdq{}}
    \PY{k}{assert} \PY{n}{np}\PY{o}{.}\PY{n}{allclose}\PY{p}{(}\PY{n}{gradients}\PY{p}{[}\PY{l+s+s1}{\PYZsq{}}\PY{l+s+s1}{dWaa}\PY{l+s+s1}{\PYZsq{}}\PY{p}{]}\PY{p}{[}\PY{l+m+mi}{1}\PY{p}{,} \PY{l+m+mi}{2}\PY{p}{]}\PY{p}{,} \PY{l+m+mf}{0.1947093}\PY{p}{)}\PY{p}{,} \PY{l+s+s2}{\PYZdq{}}\PY{l+s+s2}{Unexpected gradients. Check the rnn\PYZus{}backward call}\PY{l+s+s2}{\PYZdq{}}
    \PY{k}{assert} \PY{n}{np}\PY{o}{.}\PY{n}{allclose}\PY{p}{(}\PY{n}{gradients}\PY{p}{[}\PY{l+s+s1}{\PYZsq{}}\PY{l+s+s1}{dWya}\PY{l+s+s1}{\PYZsq{}}\PY{p}{]}\PY{p}{[}\PY{l+m+mi}{1}\PY{p}{,} \PY{l+m+mi}{2}\PY{p}{]}\PY{p}{,} \PY{o}{\PYZhy{}}\PY{l+m+mf}{0.007773876}\PY{p}{)}\PY{p}{,} \PY{l+s+s2}{\PYZdq{}}\PY{l+s+s2}{Unexpected gradients. Check the rnn\PYZus{}backward call}\PY{l+s+s2}{\PYZdq{}}
    \PY{k}{assert} \PY{o+ow}{not} \PY{n}{np}\PY{o}{.}\PY{n}{allclose}\PY{p}{(}\PY{n}{parameters}\PY{p}{[}\PY{l+s+s1}{\PYZsq{}}\PY{l+s+s1}{Wya}\PY{l+s+s1}{\PYZsq{}}\PY{p}{]}\PY{p}{,} \PY{n}{old\PYZus{}parameters}\PY{p}{[}\PY{l+s+s1}{\PYZsq{}}\PY{l+s+s1}{Wya}\PY{l+s+s1}{\PYZsq{}}\PY{p}{]}\PY{p}{)}\PY{p}{,} \PY{l+s+s2}{\PYZdq{}}\PY{l+s+s2}{parameters were not updated}\PY{l+s+s2}{\PYZdq{}}
    
    \PY{n+nb}{print}\PY{p}{(}\PY{l+s+s2}{\PYZdq{}}\PY{l+s+se}{\PYZbs{}033}\PY{l+s+s2}{[92mAll tests passed!}\PY{l+s+s2}{\PYZdq{}}\PY{p}{)}

\PY{n}{optimize\PYZus{}test}\PY{p}{(}\PY{n}{optimize}\PY{p}{)}
\end{Verbatim}
\end{tcolorbox}

    \begin{Verbatim}[commandchars=\\\{\}]
Loss = 126.50397572165389
gradients["dWaa"][1][2] = 0.19470931534713587
np.argmax(gradients["dWax"]) = 93
gradients["dWya"][1][2] = -0.007773876032002162
gradients["db"][4] = [-0.06809825]
gradients["dby"][1] = [0.01538192]
a\_last[4] = [-1.]
\textcolor{ansi-green-intense}{All tests passed!}
    \end{Verbatim}

    \textbf{Expected output}

\begin{verbatim}
Loss = 126.50397572165389
gradients["dWaa"][1][2] = 0.19470931534713587
np.argmax(gradients["dWax"]) = 93
gradients["dWya"][1][2] = -0.007773876032002162
gradients["db"][4] = [-0.06809825]
gradients["dby"][1] = [0.01538192]
a_last[4] = [-1.]
\end{verbatim}

    \#\#\# 3.2 - Training the Model

    \begin{itemize}
\tightlist
\item
  Given the dataset of dinosaur names, you'll use each line of the
  dataset (one name) as one training example.
\item
  Every 2000 steps of stochastic gradient descent, you will sample
  several randomly chosen names to see how the algorithm is doing.
\end{itemize}

\#\#\# Exercise 4 - model

Implement \texttt{model()}.

When \texttt{examples{[}index{]}} contains one dinosaur name (string),
to create an example (X, Y), you can use this:

\hypertarget{set-the-index-idx-into-the-list-of-examples}{%
\subparagraph{\texorpdfstring{Set the index \texttt{idx} into the list
of
examples}{Set the index idx into the list of examples}}\label{set-the-index-idx-into-the-list-of-examples}}

\begin{itemize}
\tightlist
\item
  Using the for-loop, walk through the shuffled list of dinosaur names
  in the list ``examples.''
\item
  For example, if there are n\_e examples, and the for-loop increments
  the index to n\_e onwards, think of how you would make the index cycle
  back to 0, so that you can continue feeding the examples into the
  model when j is n\_e, n\_e + 1, etc.
\item
  Hint: (n\_e + 1) \% n\_e equals 1, which is otherwise the `remainder'
  you get when you divide (n\_e + 1) by n\_e.
\item
  \texttt{\%} is the
  \href{https://www.freecodecamp.org/news/the-python-modulo-operator-what-does-the-symbol-mean-in-python-solved/}{modulo
  operator in python}.
\end{itemize}

\hypertarget{extract-a-single-example-from-the-list-of-examples}{%
\subparagraph{Extract a single example from the list of
examples}\label{extract-a-single-example-from-the-list-of-examples}}

\begin{itemize}
\tightlist
\item
  \texttt{single\_example}: use the \texttt{idx} index that you set
  previously to get one word from the list of examples.
\end{itemize}

    \hypertarget{convert-a-string-into-a-list-of-characters-single_example_chars}{%
\subparagraph{\texorpdfstring{Convert a string into a list of
characters:
\texttt{single\_example\_chars}}{Convert a string into a list of characters: single\_example\_chars}}\label{convert-a-string-into-a-list-of-characters-single_example_chars}}

\begin{itemize}
\tightlist
\item
  \texttt{single\_example\_chars}: A string is a list of characters.
\item
  You can use a list comprehension (recommended over for-loops) to
  generate a list of characters.
\end{itemize}

\begin{Shaded}
\begin{Highlighting}[]
\BuiltInTok{str} \OperatorTok{=} \StringTok{\textquotesingle{}I love learning\textquotesingle{}}
\NormalTok{list\_of\_chars }\OperatorTok{=}\NormalTok{ [c }\ControlFlowTok{for}\NormalTok{ c }\KeywordTok{in} \BuiltInTok{str}\NormalTok{]}
\BuiltInTok{print}\NormalTok{(list\_of\_chars)}
\end{Highlighting}
\end{Shaded}

\begin{verbatim}
['I', ' ', 'l', 'o', 'v', 'e', ' ', 'l', 'e', 'a', 'r', 'n', 'i', 'n', 'g']
\end{verbatim}

\begin{itemize}
\tightlist
\item
  For more on
  \href{https://docs.python.org/3/tutorial/datastructures.html\#list-comprehensions}{list
  comprehensions}:
\end{itemize}

    \hypertarget{convert-list-of-characters-to-a-list-of-integers-single_example_ix}{%
\subparagraph{\texorpdfstring{Convert list of characters to a list of
integers:
\texttt{single\_example\_ix}}{Convert list of characters to a list of integers: single\_example\_ix}}\label{convert-list-of-characters-to-a-list-of-integers-single_example_ix}}

\begin{itemize}
\tightlist
\item
  Create a list that contains the index numbers associated with each
  character.
\item
  Use the dictionary \texttt{char\_to\_ix}
\item
  You can combine this with the list comprehension that is used to get a
  list of characters from a string.
\end{itemize}

    \hypertarget{create-the-list-of-input-characters-x}{%
\subparagraph{\texorpdfstring{Create the list of input characters:
\texttt{X}}{Create the list of input characters: X}}\label{create-the-list-of-input-characters-x}}

\begin{itemize}
\tightlist
\item
  \texttt{rnn\_forward} uses the \textbf{\texttt{None}} value as a flag
  to set the input vector as a zero-vector.
\item
  Prepend the list {[}\textbf{\texttt{None}}{]} in front of the list of
  input character \emph{indices}.
\item
  There is more than one way to prepend a value to a list. One way is to
  add two lists together:
  \texttt{{[}\textquotesingle{}a\textquotesingle{}{]}\ +\ {[}\textquotesingle{}b\textquotesingle{}{]}}
\end{itemize}

    \hypertarget{get-the-integer-representation-of-the-newline-character-ix_newline}{%
\subparagraph{\texorpdfstring{Get the integer representation of the
newline character
\texttt{ix\_newline}}{Get the integer representation of the newline character ix\_newline}}\label{get-the-integer-representation-of-the-newline-character-ix_newline}}

\begin{itemize}
\tightlist
\item
  \texttt{ix\_newline}: The newline character signals the end of the
  dinosaur name.

  \begin{itemize}
  \tightlist
  \item
    Get the integer representation of the newline character
    \texttt{\textquotesingle{}\textbackslash{}n\textquotesingle{}}.
  \item
    Use \texttt{char\_to\_ix}
  \end{itemize}
\end{itemize}

    \hypertarget{set-the-list-of-labels-integer-representation-of-the-characters-y}{%
\subparagraph{\texorpdfstring{Set the list of labels (integer
representation of the characters):
\texttt{Y}}{Set the list of labels (integer representation of the characters): Y}}\label{set-the-list-of-labels-integer-representation-of-the-characters-y}}

\begin{itemize}
\tightlist
\item
  The goal is to train the RNN to predict the next letter in the name,
  so the labels are the list of characters that are one time-step ahead
  of the characters in the input \texttt{X}.

  \begin{itemize}
  \tightlist
  \item
    For example, \texttt{Y{[}0{]}} contains the same value as
    \texttt{X{[}1{]}}\\
  \end{itemize}
\item
  The RNN should predict a newline at the last letter, so add
  \texttt{ix\_newline} to the end of the labels.

  \begin{itemize}
  \tightlist
  \item
    Append the integer representation of the newline character to the
    end of \texttt{Y}.
  \item
    Note that \texttt{append} is an in-place operation.
  \item
    It might be easier for you to add two lists together.
  \end{itemize}
\end{itemize}

    \begin{tcolorbox}[breakable, size=fbox, boxrule=1pt, pad at break*=1mm,colback=cellbackground, colframe=cellborder]
\prompt{In}{incolor}{89}{\boxspacing}
\begin{Verbatim}[commandchars=\\\{\}]
\PY{c+c1}{\PYZsh{} UNQ\PYZus{}C4 (UNIQUE CELL IDENTIFIER, DO NOT EDIT)}
\PY{c+c1}{\PYZsh{} GRADED FUNCTION: model}

\PY{k}{def} \PY{n+nf}{model}\PY{p}{(}\PY{n}{data\PYZus{}x}\PY{p}{,} \PY{n}{ix\PYZus{}to\PYZus{}char}\PY{p}{,} \PY{n}{char\PYZus{}to\PYZus{}ix}\PY{p}{,} \PY{n}{num\PYZus{}iterations} \PY{o}{=} \PY{l+m+mi}{35000}\PY{p}{,} \PY{n}{n\PYZus{}a} \PY{o}{=} \PY{l+m+mi}{50}\PY{p}{,} \PY{n}{dino\PYZus{}names} \PY{o}{=} \PY{l+m+mi}{7}\PY{p}{,} \PY{n}{vocab\PYZus{}size} \PY{o}{=} \PY{l+m+mi}{27}\PY{p}{,} \PY{n}{verbose} \PY{o}{=} \PY{k+kc}{False}\PY{p}{)}\PY{p}{:}
    \PY{l+s+sd}{\PYZdq{}\PYZdq{}\PYZdq{}}
\PY{l+s+sd}{    Trains the model and generates dinosaur names. }
\PY{l+s+sd}{    }
\PY{l+s+sd}{    Arguments:}
\PY{l+s+sd}{    data\PYZus{}x \PYZhy{}\PYZhy{} text corpus, divided in words}
\PY{l+s+sd}{    ix\PYZus{}to\PYZus{}char \PYZhy{}\PYZhy{} dictionary that maps the index to a character}
\PY{l+s+sd}{    char\PYZus{}to\PYZus{}ix \PYZhy{}\PYZhy{} dictionary that maps a character to an index}
\PY{l+s+sd}{    num\PYZus{}iterations \PYZhy{}\PYZhy{} number of iterations to train the model for}
\PY{l+s+sd}{    n\PYZus{}a \PYZhy{}\PYZhy{} number of units of the RNN cell}
\PY{l+s+sd}{    dino\PYZus{}names \PYZhy{}\PYZhy{} number of dinosaur names you want to sample at each iteration. }
\PY{l+s+sd}{    vocab\PYZus{}size \PYZhy{}\PYZhy{} number of unique characters found in the text (size of the vocabulary)}
\PY{l+s+sd}{    }
\PY{l+s+sd}{    Returns:}
\PY{l+s+sd}{    parameters \PYZhy{}\PYZhy{} learned parameters}
\PY{l+s+sd}{    \PYZdq{}\PYZdq{}\PYZdq{}}
    
    \PY{c+c1}{\PYZsh{} Retrieve n\PYZus{}x and n\PYZus{}y from vocab\PYZus{}size}
    \PY{n}{n\PYZus{}x}\PY{p}{,} \PY{n}{n\PYZus{}y} \PY{o}{=} \PY{n}{vocab\PYZus{}size}\PY{p}{,} \PY{n}{vocab\PYZus{}size}
    
    \PY{c+c1}{\PYZsh{} Initialize parameters}
    \PY{n}{parameters} \PY{o}{=} \PY{n}{initialize\PYZus{}parameters}\PY{p}{(}\PY{n}{n\PYZus{}a}\PY{p}{,} \PY{n}{n\PYZus{}x}\PY{p}{,} \PY{n}{n\PYZus{}y}\PY{p}{)}
    
    \PY{c+c1}{\PYZsh{} Initialize loss (this is required because we want to smooth our loss)}
    \PY{n}{loss} \PY{o}{=} \PY{n}{get\PYZus{}initial\PYZus{}loss}\PY{p}{(}\PY{n}{vocab\PYZus{}size}\PY{p}{,} \PY{n}{dino\PYZus{}names}\PY{p}{)}
    
    \PY{c+c1}{\PYZsh{} Build list of all dinosaur names (training examples).}
    \PY{n}{examples} \PY{o}{=} \PY{p}{[}\PY{n}{x}\PY{o}{.}\PY{n}{strip}\PY{p}{(}\PY{p}{)} \PY{k}{for} \PY{n}{x} \PY{o+ow}{in} \PY{n}{data\PYZus{}x}\PY{p}{]}
    
    \PY{c+c1}{\PYZsh{} Shuffle list of all dinosaur names}
    \PY{n}{np}\PY{o}{.}\PY{n}{random}\PY{o}{.}\PY{n}{seed}\PY{p}{(}\PY{l+m+mi}{0}\PY{p}{)}
    \PY{n}{np}\PY{o}{.}\PY{n}{random}\PY{o}{.}\PY{n}{shuffle}\PY{p}{(}\PY{n}{examples}\PY{p}{)}
    
    \PY{c+c1}{\PYZsh{} Initialize the hidden state of your LSTM}
    \PY{n}{a\PYZus{}prev} \PY{o}{=} \PY{n}{np}\PY{o}{.}\PY{n}{zeros}\PY{p}{(}\PY{p}{(}\PY{n}{n\PYZus{}a}\PY{p}{,} \PY{l+m+mi}{1}\PY{p}{)}\PY{p}{)}
    
    \PY{c+c1}{\PYZsh{} for grading purposes}
    \PY{n}{last\PYZus{}dino\PYZus{}name} \PY{o}{=} \PY{l+s+s2}{\PYZdq{}}\PY{l+s+s2}{abc}\PY{l+s+s2}{\PYZdq{}}
    
    \PY{c+c1}{\PYZsh{} Optimization loop}
    \PY{k}{for} \PY{n}{j} \PY{o+ow}{in} \PY{n+nb}{range}\PY{p}{(}\PY{n}{num\PYZus{}iterations}\PY{p}{)}\PY{p}{:}
        
        \PY{c+c1}{\PYZsh{}\PYZsh{}\PYZsh{} START CODE HERE \PYZsh{}\PYZsh{}\PYZsh{}}
        
        \PY{c+c1}{\PYZsh{} Set the index `idx` (see instructions above)}
        \PY{n}{idx} \PY{o}{=} \PY{n}{random}\PY{o}{.}\PY{n}{randint}\PY{p}{(}\PY{l+m+mi}{0}\PY{p}{,} \PY{n}{vocab\PYZus{}size}\PY{p}{)}
        
        \PY{c+c1}{\PYZsh{} Set the input X (see instructions above)}
        \PY{n}{single\PYZus{}example} \PY{o}{=} \PY{n}{examples}\PY{p}{[}\PY{n}{idx}\PY{p}{]}
        \PY{n}{single\PYZus{}example\PYZus{}chars} \PY{o}{=} \PY{p}{[}\PY{n}{c} \PY{k}{for} \PY{n}{c} \PY{o+ow}{in} \PY{n}{single\PYZus{}example}\PY{p}{]}
        \PY{n}{single\PYZus{}example\PYZus{}ix} \PY{o}{=} \PY{p}{[}\PY{n}{char\PYZus{}to\PYZus{}ix}\PY{p}{[}\PY{n}{c}\PY{p}{]} \PY{k}{for} \PY{n}{c} \PY{o+ow}{in} \PY{n}{single\PYZus{}example\PYZus{}chars}\PY{p}{]}
        \PY{n}{X} \PY{o}{=} \PY{p}{[}\PY{k+kc}{None}\PY{p}{]} \PY{o}{+} \PY{n}{single\PYZus{}example\PYZus{}ix}
        
        \PY{c+c1}{\PYZsh{} Set the labels Y (see instructions above)}
        \PY{n}{ix\PYZus{}newline} \PY{o}{=} \PY{n}{char\PYZus{}to\PYZus{}ix}\PY{p}{[}\PY{l+s+s1}{\PYZsq{}}\PY{l+s+se}{\PYZbs{}n}\PY{l+s+s1}{\PYZsq{}}\PY{p}{]}
        \PY{n}{Y} \PY{o}{=} \PY{n}{single\PYZus{}example\PYZus{}ix} \PY{o}{+} \PY{p}{[}\PY{n}{ix\PYZus{}newline}\PY{p}{]}
        \PY{c+c1}{\PYZsh{} Perform one optimization step: Forward\PYZhy{}prop \PYZhy{}\PYZgt{} Backward\PYZhy{}prop \PYZhy{}\PYZgt{} Clip \PYZhy{}\PYZgt{} Update parameters}
        \PY{c+c1}{\PYZsh{} Choose a learning rate of 0.01}
        \PY{n}{curr\PYZus{}loss}\PY{p}{,} \PY{n}{gradients}\PY{p}{,} \PY{n}{a\PYZus{}prev} \PY{o}{=} \PY{n}{optimize}\PY{p}{(}\PY{n}{X}\PY{p}{,} \PY{n}{Y}\PY{p}{,} \PY{n}{a\PYZus{}prev}\PY{p}{,} \PY{n}{parameters}\PY{p}{,} \PY{n}{learning\PYZus{}rate} \PY{o}{=} \PY{l+m+mf}{0.01}\PY{p}{)}
        
        \PY{c+c1}{\PYZsh{}\PYZsh{}\PYZsh{} END CODE HERE \PYZsh{}\PYZsh{}\PYZsh{}}
        
        \PY{c+c1}{\PYZsh{} debug statements to aid in correctly forming X, Y}
        \PY{k}{if} \PY{n}{verbose} \PY{o+ow}{and} \PY{n}{j} \PY{o+ow}{in} \PY{p}{[}\PY{l+m+mi}{0}\PY{p}{,} \PY{n+nb}{len}\PY{p}{(}\PY{n}{examples}\PY{p}{)} \PY{o}{\PYZhy{}}\PY{l+m+mi}{1}\PY{p}{,} \PY{n+nb}{len}\PY{p}{(}\PY{n}{examples}\PY{p}{)}\PY{p}{]}\PY{p}{:}
            \PY{n+nb}{print}\PY{p}{(}\PY{l+s+s2}{\PYZdq{}}\PY{l+s+s2}{j = }\PY{l+s+s2}{\PYZdq{}} \PY{p}{,} \PY{n}{j}\PY{p}{,} \PY{l+s+s2}{\PYZdq{}}\PY{l+s+s2}{idx = }\PY{l+s+s2}{\PYZdq{}}\PY{p}{,} \PY{n}{idx}\PY{p}{,}\PY{p}{)} 
        \PY{k}{if} \PY{n}{verbose} \PY{o+ow}{and} \PY{n}{j} \PY{o+ow}{in} \PY{p}{[}\PY{l+m+mi}{0}\PY{p}{]}\PY{p}{:}
            \PY{n+nb}{print}\PY{p}{(}\PY{l+s+s2}{\PYZdq{}}\PY{l+s+s2}{single\PYZus{}example =}\PY{l+s+s2}{\PYZdq{}}\PY{p}{,} \PY{n}{single\PYZus{}example}\PY{p}{)}
            \PY{n+nb}{print}\PY{p}{(}\PY{l+s+s2}{\PYZdq{}}\PY{l+s+s2}{single\PYZus{}example\PYZus{}chars}\PY{l+s+s2}{\PYZdq{}}\PY{p}{,} \PY{n}{single\PYZus{}example\PYZus{}chars}\PY{p}{)}
            \PY{n+nb}{print}\PY{p}{(}\PY{l+s+s2}{\PYZdq{}}\PY{l+s+s2}{single\PYZus{}example\PYZus{}ix}\PY{l+s+s2}{\PYZdq{}}\PY{p}{,} \PY{n}{single\PYZus{}example\PYZus{}ix}\PY{p}{)}
            \PY{n+nb}{print}\PY{p}{(}\PY{l+s+s2}{\PYZdq{}}\PY{l+s+s2}{ X = }\PY{l+s+s2}{\PYZdq{}}\PY{p}{,} \PY{n}{X}\PY{p}{,} \PY{l+s+s2}{\PYZdq{}}\PY{l+s+se}{\PYZbs{}n}\PY{l+s+s2}{\PYZdq{}}\PY{p}{,} \PY{l+s+s2}{\PYZdq{}}\PY{l+s+s2}{Y =       }\PY{l+s+s2}{\PYZdq{}}\PY{p}{,} \PY{n}{Y}\PY{p}{,} \PY{l+s+s2}{\PYZdq{}}\PY{l+s+se}{\PYZbs{}n}\PY{l+s+s2}{\PYZdq{}}\PY{p}{)}
        
        \PY{c+c1}{\PYZsh{} to keep the loss smooth.}
        \PY{n}{loss} \PY{o}{=} \PY{n}{smooth}\PY{p}{(}\PY{n}{loss}\PY{p}{,} \PY{n}{curr\PYZus{}loss}\PY{p}{)}

        \PY{c+c1}{\PYZsh{} Every 2000 Iteration, generate \PYZdq{}n\PYZdq{} characters thanks to sample() to check if the model is learning properly}
        \PY{k}{if} \PY{n}{j} \PY{o}{\PYZpc{}} \PY{l+m+mi}{2000} \PY{o}{==} \PY{l+m+mi}{0}\PY{p}{:}
            
            \PY{n+nb}{print}\PY{p}{(}\PY{l+s+s1}{\PYZsq{}}\PY{l+s+s1}{Iteration: }\PY{l+s+si}{\PYZpc{}d}\PY{l+s+s1}{, Loss: }\PY{l+s+si}{\PYZpc{}f}\PY{l+s+s1}{\PYZsq{}} \PY{o}{\PYZpc{}} \PY{p}{(}\PY{n}{j}\PY{p}{,} \PY{n}{loss}\PY{p}{)} \PY{o}{+} \PY{l+s+s1}{\PYZsq{}}\PY{l+s+se}{\PYZbs{}n}\PY{l+s+s1}{\PYZsq{}}\PY{p}{)}
            
            \PY{c+c1}{\PYZsh{} The number of dinosaur names to print}
            \PY{n}{seed} \PY{o}{=} \PY{l+m+mi}{0}
            \PY{k}{for} \PY{n}{name} \PY{o+ow}{in} \PY{n+nb}{range}\PY{p}{(}\PY{n}{dino\PYZus{}names}\PY{p}{)}\PY{p}{:}
                
                \PY{c+c1}{\PYZsh{} Sample indices and print them}
                \PY{n}{sampled\PYZus{}indices} \PY{o}{=} \PY{n}{sample}\PY{p}{(}\PY{n}{parameters}\PY{p}{,} \PY{n}{char\PYZus{}to\PYZus{}ix}\PY{p}{,} \PY{n}{seed}\PY{p}{)}
                \PY{n}{last\PYZus{}dino\PYZus{}name} \PY{o}{=} \PY{n}{get\PYZus{}sample}\PY{p}{(}\PY{n}{sampled\PYZus{}indices}\PY{p}{,} \PY{n}{ix\PYZus{}to\PYZus{}char}\PY{p}{)}
                \PY{n+nb}{print}\PY{p}{(}\PY{n}{last\PYZus{}dino\PYZus{}name}\PY{o}{.}\PY{n}{replace}\PY{p}{(}\PY{l+s+s1}{\PYZsq{}}\PY{l+s+se}{\PYZbs{}n}\PY{l+s+s1}{\PYZsq{}}\PY{p}{,} \PY{l+s+s1}{\PYZsq{}}\PY{l+s+s1}{\PYZsq{}}\PY{p}{)}\PY{p}{)}
                
                \PY{n}{seed} \PY{o}{+}\PY{o}{=} \PY{l+m+mi}{1}  \PY{c+c1}{\PYZsh{} To get the same result (for grading purposes), increment the seed by one. }
      
            \PY{n+nb}{print}\PY{p}{(}\PY{l+s+s1}{\PYZsq{}}\PY{l+s+se}{\PYZbs{}n}\PY{l+s+s1}{\PYZsq{}}\PY{p}{)}
        
    \PY{k}{return} \PY{n}{parameters}\PY{p}{,} \PY{n}{last\PYZus{}dino\PYZus{}name}
\end{Verbatim}
\end{tcolorbox}

    When you run the following cell, you should observe your model
outputting random-looking characters at the first iteration. After a few
thousand iterations, your model should learn to generate
reasonable-looking names.

    \begin{tcolorbox}[breakable, size=fbox, boxrule=1pt, pad at break*=1mm,colback=cellbackground, colframe=cellborder]
\prompt{In}{incolor}{90}{\boxspacing}
\begin{Verbatim}[commandchars=\\\{\}]
\PY{n}{parameters}\PY{p}{,} \PY{n}{last\PYZus{}name} \PY{o}{=} \PY{n}{model}\PY{p}{(}\PY{n}{data}\PY{o}{.}\PY{n}{split}\PY{p}{(}\PY{l+s+s2}{\PYZdq{}}\PY{l+s+se}{\PYZbs{}n}\PY{l+s+s2}{\PYZdq{}}\PY{p}{)}\PY{p}{,} \PY{n}{ix\PYZus{}to\PYZus{}char}\PY{p}{,} \PY{n}{char\PYZus{}to\PYZus{}ix}\PY{p}{,} \PY{l+m+mi}{22001}\PY{p}{,} \PY{n}{verbose} \PY{o}{=} \PY{k+kc}{True}\PY{p}{)}

\PY{k}{assert} \PY{n}{last\PYZus{}name} \PY{o}{==} \PY{l+s+s1}{\PYZsq{}}\PY{l+s+s1}{Trodonosaurus}\PY{l+s+se}{\PYZbs{}n}\PY{l+s+s1}{\PYZsq{}}\PY{p}{,} \PY{l+s+s2}{\PYZdq{}}\PY{l+s+s2}{Wrong expected output}\PY{l+s+s2}{\PYZdq{}}
\PY{n+nb}{print}\PY{p}{(}\PY{l+s+s2}{\PYZdq{}}\PY{l+s+se}{\PYZbs{}033}\PY{l+s+s2}{[92mAll tests passed!}\PY{l+s+s2}{\PYZdq{}}\PY{p}{)}
\end{Verbatim}
\end{tcolorbox}

    \begin{Verbatim}[commandchars=\\\{\}]
j =  0 idx =  2
single\_example = ilokelesia
single\_example\_chars ['i', 'l', 'o', 'k', 'e', 'l', 'e', 's', 'i', 'a']
single\_example\_ix [9, 12, 15, 11, 5, 12, 5, 19, 9, 1]
 X =  [None, 9, 12, 15, 11, 5, 12, 5, 19, 9, 1]
 Y =        [9, 12, 15, 11, 5, 12, 5, 19, 9, 1, 0]

Iteration: 0, Loss: 23.084044

Nkzxwtdmfqoeyhsqwasjjjvu
Kneb
Kzxwtdmfqoeyhsqwasjjjvu
Neb
Zxwtdmfqoeyhsqwasjjjvu
Eb
Xwtdmfqoeyhsqwasjjjvu


j =  1535 idx =  11
j =  1536 idx =  22
Iteration: 2000, Loss: 23.585376

Nesotomaauposaurur
Lokachus
Lusiparatocxelus
Naccherocautosaurus
Vrhongheritosaurus
Deaerus
Tosaurus


Iteration: 4000, Loss: 10.708640

Horshamosaurus
Chebndosekanosaurus
Coelosaurus
Hor
Torsaugosaurus
Chaburacertoscherocauchaighusicoelurus
Torcherocomauris


Iteration: 6000, Loss: 5.266862

Horshamosaurus
Chubothaenhamdhubotortyrus
Coelosaurus
Horbhoposaurus
Torthenotos
Chaburacdus
Torcherocoptos


Iteration: 8000, Loss: 20.399887

Lirytsaurusaurushgros
Linebatopaantosaurus
Lusus
Lenamss
Tusif
Degnsaulanusil
Sisanos


Iteration: 10000, Loss: 27.867737

Llttus
Loga
Lusrurumus
Lla
Tostheltogteusubuoleso
Hh
Sus


Iteration: 12000, Loss: 29.108458

Lhvorthncnoetlnos
Hog
Itosigjllolugorpdolelus
Lac
Tos
Daa
Rus


Iteration: 14000, Loss: 29.445048

Losus
Loraahrugaaoshuoghurecvairra
Lusus
Liaaesur
Trruahamoiterruathaerud
Heaerrga
Rusaprus


Iteration: 16000, Loss: 30.505131

Letttomanos
Lhhaaesoracmustherus
Lusiupsaurnursus
Lac
Usos
Daacusacanrsupesiplaxchos
Rus


Iteration: 18000, Loss: 27.992077

Litutocianraurustesamhus
Lor
Lusauopaurnunpsoaurinos
La
Tus
Ha
Rus


Iteration: 20000, Loss: 28.856063

Lotrsicelor
Hpra
Hytrr
Lea
Tus
Ac
Rus


Iteration: 22000, Loss: 29.006625

Lavusaropropthprtatilhrq
Lea
Lyutumeerhgthoseatikitucosavapaupcteglaphaconusaur
Laa
Tutumehoravhoseatikitucosavapaupaudhgconadontphori
Da
Rtsasauritnorubrimksq


    \end{Verbatim}

    \begin{Verbatim}[commandchars=\\\{\}]

        ---------------------------------------------------------------------------

        AssertionError                            Traceback (most recent call last)

        <ipython-input-90-725c093d6b91> in <module>
          1 parameters, last\_name = model(data.split("\textbackslash{}n"), ix\_to\_char, char\_to\_ix, 22001, verbose = True)
          2 
    ----> 3 assert last\_name == 'Trodonosaurus\textbackslash{}n', "Wrong expected output"
          4 print("\textbackslash{}033[92mAll tests passed!")


        AssertionError: Wrong expected output

    \end{Verbatim}

    \textbf{Expected output}

\begin{verbatim}
...
Iteration: 22000, Loss: 22.728886

Onustreofkelus
Llecagosaurus
Mystolojmiaterltasaurus
Ola
Yuskeolongus
Eiacosaurus
Trodonosaurus
\end{verbatim}

    \hypertarget{conclusion}{%
\subsubsection{Conclusion}\label{conclusion}}

You can see that your algorithm has started to generate plausible
dinosaur names towards the end of training. At first, it was generating
random characters, but towards the end you could begin to see dinosaur
names with cool endings. Feel free to run the algorithm even longer and
play with hyperparameters to see if you can get even better results! Our
implementation generated some really cool names like \texttt{maconucon},
\texttt{marloralus} and \texttt{macingsersaurus}. Your model hopefully
also learned that dinosaur names tend to end in \texttt{saurus},
\texttt{don}, \texttt{aura}, \texttt{tor}, etc.

If your model generates some non-cool names, don't blame the model
entirely -- not all actual dinosaur names sound cool. (For example,
\texttt{dromaeosauroides} is an actual dinosaur name and is in the
training set.) But this model should give you a set of candidates from
which you can pick the coolest!

This assignment used a relatively small dataset, so that you're able to
train an RNN quickly on a CPU. Training a model of the English language
requires a much bigger dataset, and usually much more computation, and
could run for many hours on GPUs. We ran our dinosaur name for quite
some time, and so far our favorite name is the great, the fierce, the
undefeated: \textbf{Mangosaurus}!

    \hypertarget{congratulations}{%
\subsubsection{Congratulations!}\label{congratulations}}

You've finished the graded portion of this notebook and created a
working language model! Awesome job.

By now, you've:

\begin{itemize}
\tightlist
\item
  Stored text data for processing using an RNN
\item
  Built a character-level text generation model
\item
  Explored the vanishing/exploding gradient problem in RNNs
\item
  Applied gradient clipping to avoid exploding gradients
\end{itemize}

You've also hopefully generated some dinosaur names that are cool enough
to please you and also avoid the wrath of the dinosaurs. If you had fun
with the assignment, be sure not to miss the ungraded portion, where
you'll be able to generate poetry like the Bard Himself. Good luck and
have fun!

    \#\# 4 - Writing like Shakespeare (OPTIONAL/UNGRADED)

The rest of this notebook is optional and is not graded, but it's quite
fun and informative, so you're highly encouraged to try it out!

A similar task to character-level text generation (but more complicated)
is generating Shakespearean poems. Instead of learning from a dataset of
dinosaur names, you can use a collection of Shakespearean poems. Using
LSTM cells, you can learn longer-term dependencies that span many
characters in the text--e.g., where a character appearing somewhere a
sequence can influence what should be a different character, much later
in the sequence. These long-term dependencies were less important with
dinosaur names, since the names were quite short.

Let's become poets!

Below, you can implement a Shakespeare poem generator with Keras. Run
the following cell to load the required packages and models. This may
take a few minutes.

    \begin{tcolorbox}[breakable, size=fbox, boxrule=1pt, pad at break*=1mm,colback=cellbackground, colframe=cellborder]
\prompt{In}{incolor}{ }{\boxspacing}
\begin{Verbatim}[commandchars=\\\{\}]
\PY{k+kn}{from} \PY{n+nn}{\PYZus{}\PYZus{}future\PYZus{}\PYZus{}} \PY{k+kn}{import} \PY{n}{print\PYZus{}function}
\PY{k+kn}{from} \PY{n+nn}{tensorflow}\PY{n+nn}{.}\PY{n+nn}{keras}\PY{n+nn}{.}\PY{n+nn}{callbacks} \PY{k+kn}{import} \PY{n}{LambdaCallback}
\PY{k+kn}{from} \PY{n+nn}{tensorflow}\PY{n+nn}{.}\PY{n+nn}{keras}\PY{n+nn}{.}\PY{n+nn}{models} \PY{k+kn}{import} \PY{n}{Model}\PY{p}{,} \PY{n}{load\PYZus{}model}\PY{p}{,} \PY{n}{Sequential}
\PY{k+kn}{from} \PY{n+nn}{tensorflow}\PY{n+nn}{.}\PY{n+nn}{keras}\PY{n+nn}{.}\PY{n+nn}{layers} \PY{k+kn}{import} \PY{n}{Dense}\PY{p}{,} \PY{n}{Activation}\PY{p}{,} \PY{n}{Dropout}\PY{p}{,} \PY{n}{Input}\PY{p}{,} \PY{n}{Masking}
\PY{k+kn}{from} \PY{n+nn}{tensorflow}\PY{n+nn}{.}\PY{n+nn}{keras}\PY{n+nn}{.}\PY{n+nn}{layers} \PY{k+kn}{import} \PY{n}{LSTM}
\PY{k+kn}{from} \PY{n+nn}{tensorflow}\PY{n+nn}{.}\PY{n+nn}{keras}\PY{n+nn}{.}\PY{n+nn}{utils} \PY{k+kn}{import} \PY{n}{get\PYZus{}file}
\PY{k+kn}{from} \PY{n+nn}{tensorflow}\PY{n+nn}{.}\PY{n+nn}{keras}\PY{n+nn}{.}\PY{n+nn}{preprocessing}\PY{n+nn}{.}\PY{n+nn}{sequence} \PY{k+kn}{import} \PY{n}{pad\PYZus{}sequences}
\PY{k+kn}{from} \PY{n+nn}{shakespeare\PYZus{}utils} \PY{k+kn}{import} \PY{o}{*}
\PY{k+kn}{import} \PY{n+nn}{sys}
\PY{k+kn}{import} \PY{n+nn}{io}
\end{Verbatim}
\end{tcolorbox}

    To save you some time, a model has already been trained for
\textasciitilde1000 epochs on a collection of Shakespearean poems called
``\href{shakespeare.txt}{The Sonnets}.''

    Let's train the model for one more epoch. When it finishes training for
an epoch (this will also take a few minutes), you can run
\texttt{generate\_output}, which will prompt you for an input
(\texttt{\textless{}}40 characters). The poem will start with your
sentence, and your RNN Shakespeare will complete the rest of the poem
for you! For example, try, ``Forsooth this maketh no sense'' (without
the quotation marks!). Depending on whether you include the space at the
end, your results might also differ, so try it both ways, and try other
inputs as well.

    \begin{tcolorbox}[breakable, size=fbox, boxrule=1pt, pad at break*=1mm,colback=cellbackground, colframe=cellborder]
\prompt{In}{incolor}{ }{\boxspacing}
\begin{Verbatim}[commandchars=\\\{\}]
\PY{n}{print\PYZus{}callback} \PY{o}{=} \PY{n}{LambdaCallback}\PY{p}{(}\PY{n}{on\PYZus{}epoch\PYZus{}end}\PY{o}{=}\PY{n}{on\PYZus{}epoch\PYZus{}end}\PY{p}{)}

\PY{n}{model}\PY{o}{.}\PY{n}{fit}\PY{p}{(}\PY{n}{x}\PY{p}{,} \PY{n}{y}\PY{p}{,} \PY{n}{batch\PYZus{}size}\PY{o}{=}\PY{l+m+mi}{128}\PY{p}{,} \PY{n}{epochs}\PY{o}{=}\PY{l+m+mi}{1}\PY{p}{,} \PY{n}{callbacks}\PY{o}{=}\PY{p}{[}\PY{n}{print\PYZus{}callback}\PY{p}{]}\PY{p}{)}
\end{Verbatim}
\end{tcolorbox}

    \begin{tcolorbox}[breakable, size=fbox, boxrule=1pt, pad at break*=1mm,colback=cellbackground, colframe=cellborder]
\prompt{In}{incolor}{ }{\boxspacing}
\begin{Verbatim}[commandchars=\\\{\}]
\PY{c+c1}{\PYZsh{} Run this cell to try with different inputs without having to re\PYZhy{}train the model }
\PY{n}{generate\PYZus{}output}\PY{p}{(}\PY{p}{)}
\end{Verbatim}
\end{tcolorbox}

    \hypertarget{congratulations-on-finishing-this-notebook}{%
\subsubsection{Congratulations on finishing this
notebook!}\label{congratulations-on-finishing-this-notebook}}

The RNN Shakespeare model is very similar to the one you built for
dinosaur names. The only major differences are: - LSTMs instead of the
basic RNN to capture longer-range dependencies - The model is a deeper,
stacked LSTM model (2 layer) - Using Keras instead of Python to simplify
the code

    \#\# 5 - References - This exercise took inspiration from Andrej
Karpathy's implementation:
https://gist.github.com/karpathy/d4dee566867f8291f086. To learn more
about text generation, also check out Karpathy's
\href{http://karpathy.github.io/2015/05/21/rnn-effectiveness/}{blog
post}.

    \begin{tcolorbox}[breakable, size=fbox, boxrule=1pt, pad at break*=1mm,colback=cellbackground, colframe=cellborder]
\prompt{In}{incolor}{ }{\boxspacing}
\begin{Verbatim}[commandchars=\\\{\}]

\end{Verbatim}
\end{tcolorbox}


    % Add a bibliography block to the postdoc
    
    
    
\end{document}
